% Pertemuan 14 Praktikum Pemrograman (IF25-11002). Dihasilkan oleh tools/build.py dari tools/konten/p14.py.
% Kompilasi: xelatex P14.tex (dua kali). Catatan: xelatex -jobname=P14-catatan "\def\catatan{1}% Pertemuan 14 Praktikum Pemrograman (IF25-11002). Dihasilkan oleh tools/build.py dari tools/konten/p14.py.
% Kompilasi: xelatex P14.tex (dua kali). Catatan: xelatex -jobname=P14-catatan "\def\catatan{1}% Pertemuan 14 Praktikum Pemrograman (IF25-11002). Dihasilkan oleh tools/build.py dari tools/konten/p14.py.
% Kompilasi: xelatex P14.tex (dua kali). Catatan: xelatex -jobname=P14-catatan "\def\catatan{1}% Pertemuan 14 Praktikum Pemrograman (IF25-11002). Dihasilkan oleh tools/build.py dari tools/konten/p14.py.
% Kompilasi: xelatex P14.tex (dua kali). Catatan: xelatex -jobname=P14-catatan "\def\catatan{1}\input{P14.tex}"
\documentclass[t]{beamer}
\usetheme{ITERA}
\input{catatan-preamble.tex}
\renewcommand\ITERAmeeting{Pertemuan 14}
\begin{document}
\SlideJudul{IF25-11002 PRAKTIKUM PEMROGRAMAN}{Pertemuan 14\\Rekursi dan Struct}{Fungsi yang memanggil dirinya sendiri, dan tipe data buatan sendiri untuk mengelompokkan data}{Tim Pengajar}{Tahun 2026. Sejajar dengan Pertemuan 14 Algoritma Pemrograman: Rekursi dan Struct.}
\note{Buka dengan menanyakan satu hal dari pertemuan lalu yang masih membingungkan.}

\begin{frame}{Dua mata kuliah, satu minggu}
\framesubtitle{Sebelum mulai}
Topik minggu ini sama di dua mata kuliah, dengan peran yang berbeda.
\vspace{10pt}

\begin{columns}[T]
\begin{column}{0.47\textwidth}
\begin{Blok}{Algoritma: Rekursi dan Struct}
Definisi rekursif, kasus basis, pohon rekursi, serta record atau tipe bentukan sebagai kumpulan field.
\end{Blok}
\end{column}
\begin{column}{0.47\textwidth}
\begin{BlokGelap}{Praktikum: Rekursi dan struct}
Menulis fungsi rekursif dengan kasus basis yang tepat, mendefinisikan \texttt{struct}, dan memproses array of struct.
\end{BlokGelap}
\end{column}
\end{columns}
\note{Tanyakan apa yang sudah dibahas di Algoritma minggu ini. Gunakan jawabannya sebagai jembatan ke praktikum.}
\end{frame}

\begin{frame}{Saat pulang nanti, kalian bisa}
\framesubtitle{Target hari ini}
\begin{columns}[T]
\begin{column}{0.47\textwidth}
\begin{Langkah}{1}{Menerapkan}
Menulis fungsi rekursif dengan kasus basis dan langkah rekursif yang tepat, serta mendefinisikan \texttt{struct} dan memproses array of struct

{\footnotesize\color{ITERAInkMuted}Sub-CPMK 14.1, CPMK0607}
\end{Langkah}
\end{column}
\begin{column}{0.47\textwidth}
\begin{Langkah}{2}{Menjelaskan}
Menelusuri pemanggilan rekursif dengan pohon pemanggilan atau tumpukan panggilan, dan menjelaskan peran kasus basis serta akses anggota struct

{\footnotesize\color{ITERAInkMuted}Sub-CPMK 14.2, CPMK0608}
\end{Langkah}
\end{column}
\end{columns}
\vspace{8pt}
\begin{Catatan}
Dua kemampuan ini dinilai sama berat di Exit Ticket, tugas, UTS, dan UAS.
\end{Catatan}
\note{Bacakan target dengan pelan. Kata kuncinya menerapkan dan menjelaskan.}
\end{frame}

\begin{frame}{Ingat minggu lalu}
\framesubtitle{Review, 5 menit}
\begin{columns}[T]
\begin{column}{0.31\textwidth}
\begin{BlokGelap}{Pertanyaan 1}
Apa beda pass by value dan pass by reference?
\end{BlokGelap}
\end{column}
\begin{column}{0.31\textwidth}
\begin{BlokGelap}{Pertanyaan 2}
Mengapa fungsi non-void harus punya \texttt{return} di semua jalur?
\end{BlokGelap}
\end{column}
\begin{column}{0.31\textwidth}
\begin{BlokGelap}{Pertanyaan 3}
Apakah array di dalam fungsi adalah salinan?
\end{BlokGelap}
\end{column}
\end{columns}
\vspace{8pt}
Jawab di kertas dulu, lalu cocokkan dengan teman sebelah.\note{Pilih dua mahasiswa untuk menjawab. Koreksi singkat, jangan mengajar ulang.}
\end{frame}

\Bagian{1}{Masalah pembuka}{Data buku perpustakaan}
\note{Masalah pembuka 10 menit.}

\begin{frame}[fragile]{Data buku perpustakaan}
\framesubtitle{Masalah minggu ini}
\begin{columns}[T]
\begin{column}{0.55\textwidth}
{\small Perpustakaan prodi mencatat setiap buku dengan kode, judul, penulis, tahun, dan stok. Dengan materi sejauh ini, kalian butuh lima array sejajar dan harus menjaga indeksnya selalu sama.

\vspace{6pt}
Selain itu, beberapa pertanyaan lebih mudah dirumuskan dari ukuran yang lebih kecil, misalnya: total stok n buku adalah stok buku terakhir ditambah total stok n − 1 buku sebelumnya.\par}
\vspace{6pt}
\begin{Catatan}
\textbf{Diskusi 2 menit.} Bagaimana menyimpan lima keterangan satu buku dalam satu variabel? Bagaimana menghitung total stok tanpa perulangan?
\end{Catatan}
\end{column}
\begin{column}{0.41\textwidth}
\begin{Keluaran}[]
[B01] Algoritma (2019) stok 5
[B02] C++ Dasar (2021) stok 8
[B03] Struktur Data (2018) stok 10
Total stok: 23
Buku terlama: Struktur Data
\end{Keluaran}
\end{column}
\end{columns}
\note{Tampung beberapa jawaban tanpa menilai benar salah.}
\end{frame}

\begin{frame}{Dari langkah ke instruksi}
\framesubtitle{Sambungan dengan Algoritma}
\begin{columns}[T]
\begin{column}{0.31\textwidth}
\begin{Langkah}{1}{Kelompokkan}
Satu buku punya beberapa field: kode, judul, penulis, tahun, stok.
\end{Langkah}
\end{column}
\begin{column}{0.31\textwidth}
\begin{Langkah}{2}{Perkecil masalah}
Total n buku = buku terakhir + total n − 1 buku; total 0 buku = 0.
\end{Langkah}
\end{column}
\begin{column}{0.31\textwidth}
\begin{Langkah}{3}{Terjemahkan}
\texttt{struct Buku} dan fungsi rekursif dengan kasus basis.
\end{Langkah}
\end{column}
\end{columns}
\vspace{10pt}
Langkah 1 dan 2 dilatih di Algoritma. Langkah 3 kita kerjakan hari ini dengan C++.\note{Hubungkan eksplisit ke Algoritma minggu ini.}
\end{frame}

\Bagian{2}{Coba dulu, baru dijelaskan}{25 menit eksperimen di GDB Online}
\note{Struggle 25 menit. Berkeliling, jangan menjelaskan materi dulu.}

\begin{frame}{Misi kalian}
\framesubtitle{Struggle, 25 menit}
\begin{columns}[T]
\begin{column}{0.31\textwidth}
\begin{Langkah}{1}{Hitung mundur}
Tulis fungsi yang menampilkan n, n − 1, sampai 1 tanpa \texttt{for} dan \texttt{while}.
\end{Langkah}
\end{column}
\begin{column}{0.31\textwidth}
\begin{Langkah}{2}{Berhenti}
Apa yang terjadi bila fungsi itu tidak punya syarat berhenti?
\end{Langkah}
\end{column}
\begin{column}{0.31\textwidth}
\begin{Langkah}{3}{Kelompokkan}
Simpan nama, NIM, dan IPK satu mahasiswa dalam satu variabel.
\end{Langkah}
\end{column}
\end{columns}
\vspace{8pt}
\begin{columns}[T]
\begin{column}{0.47\textwidth}
\begin{Blok}{Boleh}
Bertanya ke teman, mengubah apa saja, sengaja membuat error.
\end{Blok}
\end{column}
\begin{column}{0.47\textwidth}
\begin{BlokAlert}{Belum dulu}
Mencari jawaban jadi di internet atau AI.
\end{BlokAlert}
\end{column}
\end{columns}
\note{Catat kesalahan yang sering muncul untuk dibahas di debrief.}
\end{frame}

\begin{frame}{Apa yang kalian temukan?}
\framesubtitle{Debrief struggle}
\begin{columns}[T]
\begin{column}{0.31\textwidth}
\begin{BlokGelap}{Pertanyaan 1}
Bagaimana fungsi bisa memanggil dirinya sendiri?
\end{BlokGelap}
\end{column}
\begin{column}{0.31\textwidth}
\begin{BlokGelap}{Pertanyaan 2}
Syarat apa yang membuatnya berhenti?
\end{BlokGelap}
\end{column}
\begin{column}{0.31\textwidth}
\begin{BlokGelap}{Pertanyaan 3}
Apa yang kalian coba untuk menyatukan tiga data mahasiswa?
\end{BlokGelap}
\end{column}
\end{columns}
\vspace{10pt}
Jawaban kalian akan kita uji di bagian berikutnya.\note{Tulis jawaban di papan; buktikan atau patahkan di bagian materi.}
\end{frame}

\Bagian{3}{Memanggil diri sendiri dan tipe buatan}{Kasus basis, pohon pemanggilan, struct, dan array of struct}
\note{Materi 40 menit. Tunjukkan setiap contoh langsung di GDB Online.}

\begin{frame}[fragile]{Rekursi: kasus basis dan langkah rekursif}
\framesubtitle{Faktorial}
\begin{columns}[T]
\begin{column}{0.55\textwidth}
\begin{Kode}[]{faktorial.cpp}
int faktorial(int n) {
    if (n == 0) return 1;
    return n * faktorial(n - 1);
}

int main() {
    cout << faktorial(5) << endl;
    cout << faktorial(0) << endl;
    return 0;
}
\end{Kode}
\end{column}
\begin{column}{0.41\textwidth}
\only<1>{\begin{TebakDulu}
{\ITERAKickerFont\fontsize{10pt}{12pt}\selectfont\color{ITERAAlert}TEBAK DULU}\par\vspace{4pt}
Sebelum dijalankan, tulis keluarannya di kertas.
\end{TebakDulu}
}\only<2>{\KeluaranFile[]{keluaran/P14/k001.txt}
\begin{Catatan}
Kasus basis menghentikan rekursi. Langkah rekursif harus selalu mendekati kasus basis.
\end{Catatan}
}
\end{column}
\end{columns}
\note<1>{Minta mahasiswa menebak keluaran sebelum membuka slide berikutnya. Tanyakan satu atau dua tebakan yang berbeda, lalu buktikan.}
\end{frame}

\begin{frame}{Menelusuri faktorial(4)}
\framesubtitle{Tumpukan pemanggilan}
\small
\begin{tabular}{@{}>{\raggedright\arraybackslash}p{249pt}>{\raggedright\arraybackslash}p{307pt}>{\raggedright\arraybackslash}p{307pt}@{}}
\kepala{Pemanggilan} & \kepala{Menunggu} & \kepala{Nilai kembali}\\
faktorial(4) & 4 × faktorial(3) & 4 × 6 = 24\\
\barisgenap faktorial(3) & 3 × faktorial(2) & 3 × 2 = 6\\
faktorial(2) & 2 × faktorial(1) & 2 × 1 = 2\\
\barisgenap faktorial(1) & 1 × faktorial(0) & 1 × 1 = 1\\
faktorial(0) & kasus basis & 1\\
\garisdasar
\end{tabular}
\vspace{10pt}

Pemanggilan turun sampai kasus basis, lalu nilai kembali naik dari bawah ke atas.
\end{frame}

\begin{frame}[fragile]{Rekursi bercabang dan rekursi pada array}
\framesubtitle{Fibonacci dan jumlah array}
\begin{columns}[T]
\begin{column}{0.60\textwidth}
\begin{Kode}[listing options={style=ITERACodeMini}]{fibarray.cpp}
int fib(int n) {
    if (n <= 1) return n;
    return fib(n - 1) + fib(n - 2);
}
int jumlah(int a[], int n) {
    if (n == 0) return 0;
    return a[n - 1] + jumlah(a, n - 1);
}
int main() {
    int a[4] = {5, 8, 10, 2};
    cout << fib(7) << " " << jumlah(a, 4) << endl;
    return 0;
}
\end{Kode}
\end{column}
\begin{column}{0.36\textwidth}
\only<1>{\begin{TebakDulu}
{\ITERAKickerFont\fontsize{10pt}{12pt}\selectfont\color{ITERAAlert}TEBAK DULU}\par\vspace{4pt}
Sebelum dijalankan, tulis keluarannya di kertas.
\end{TebakDulu}
}\only<2>{\KeluaranFile[listing options={style=ITERAPlainKecil}]{keluaran/P14/k002.txt}
\begin{Catatan}
\texttt{fib} memanggil dirinya dua kali, sehingga pohon pemanggilannya melebar. \texttt{jumlah} memproses elemen terakhir lalu sisanya.
\end{Catatan}
}
\end{column}
\end{columns}
\note<1>{Minta mahasiswa menebak keluaran sebelum membuka slide berikutnya. Tanyakan satu atau dua tebakan yang berbeda, lalu buktikan.}
\end{frame}

\begin{frame}[fragile]{struct: tipe data buatan sendiri}
\framesubtitle{Mengelompokkan field}
\begin{columns}[T]
\begin{column}{0.60\textwidth}
\begin{Kode}[listing options={style=ITERACodeMini}]{struct.cpp}
struct Mahasiswa {
    string nama;
    string nim;
    double ipk;
};
int main() {
    Mahasiswa m;
    m.nama = "Rani";
    m.nim = "126140001";
    m.ipk = 3.72;
    cout << m.nama << " " << m.nim << " " << m.ipk << endl;
    return 0;
}
\end{Kode}
\end{column}
\begin{column}{0.36\textwidth}
\only<1>{\begin{TebakDulu}
{\ITERAKickerFont\fontsize{10pt}{12pt}\selectfont\color{ITERAAlert}TEBAK DULU}\par\vspace{4pt}
Sebelum dijalankan, tulis keluarannya di kertas.
\end{TebakDulu}
}\only<2>{\KeluaranFile[listing options={style=ITERAPlainKecil}]{keluaran/P14/k003.txt}
\begin{Catatan}
Definisi \texttt{struct} diakhiri titik koma. Field diakses dengan titik: \texttt{m.nama}.
\end{Catatan}
}
\end{column}
\end{columns}
\note<1>{Minta mahasiswa menebak keluaran sebelum membuka slide berikutnya. Tanyakan satu atau dua tebakan yang berbeda, lalu buktikan.}
\end{frame}

\begin{frame}[fragile]{Array of struct}
\framesubtitle{Banyak data terstruktur}
\begin{columns}[T]
\begin{column}{0.60\textwidth}
\begin{Kode}[listing options={style=ITERACodeKecil}]{arraystruct.cpp}
struct Buku { string judul; int tahun; };

int main() {
    Buku rak[3] = {{"Algoritma", 2019},
                   {"C++ Dasar", 2021},
                   {"Struktur Data", 2018}};
    int iLama = 0;
    for (int i = 1; i < 3; i++)
        if (rak[i].tahun < rak[iLama].tahun) iLama = i;
    cout << rak[iLama].judul << endl;
    return 0;
}
\end{Kode}
\end{column}
\begin{column}{0.36\textwidth}
\only<1>{\begin{TebakDulu}
{\ITERAKickerFont\fontsize{10pt}{12pt}\selectfont\color{ITERAAlert}TEBAK DULU}\par\vspace{4pt}
Sebelum dijalankan, tulis keluarannya di kertas.
\end{TebakDulu}
}\only<2>{\KeluaranFile[listing options={style=ITERAPlainKecil}]{keluaran/P14/k004.txt}
\begin{Catatan}
\texttt{rak[i].tahun}: indeks dulu untuk memilih buku, lalu titik untuk memilih field.
\end{Catatan}
}
\end{column}
\end{columns}
\note<1>{Minta mahasiswa menebak keluaran sebelum membuka slide berikutnya. Tanyakan satu atau dua tebakan yang berbeda, lalu buktikan.}
\end{frame}

\begin{frame}[fragile]{Tebak keluarannya}
\framesubtitle{Latihan menjelaskan, 3 menit}
\begin{columns}[T]
\begin{column}{0.56\textwidth}
\begin{Kode}[]{tebak.cpp}
int f(int n) {
    if (n <= 0) return 0;
    if (n % 2 == 0) return n + f(n - 1);
    return f(n - 1);
}

int main() {
    cout << f(6) << " " << f(3) << endl;
    return 0;
}
\end{Kode}
\end{column}
\begin{column}{0.40\textwidth}
\begin{Catatan}
Tulis keluarannya di kertas, karakter demi karakter. Jangan dijalankan dulu.
\end{Catatan}
\end{column}
\end{columns}
\note{Contoh soal Bagian B. Beri waktu 3 menit sebelum slide jawaban.}
\end{frame}

\begin{frame}[fragile]{Tebak keluarannya: jawaban}
\framesubtitle{Latihan menjelaskan}
\begin{columns}[T]
\begin{column}{0.56\textwidth}
\begin{Kode}[]{tebak.cpp}
int f(int n) {
    if (n <= 0) return 0;
    if (n % 2 == 0) return n + f(n - 1);
    return f(n - 1);
}

int main() {
    cout << f(6) << " " << f(3) << endl;
    return 0;
}
\end{Kode}
\end{column}
\begin{column}{0.40\textwidth}
\begin{Keluaran}[]
12 2
\end{Keluaran}
\begin{Catatan}
Hanya bilangan genap yang dijumlahkan: f(6) = 6 + 4 + 2 = 12, f(3) = 2.
\end{Catatan}
\end{column}
\end{columns}
\note{Tanyakan jawaban yang paling sering salah, lalu telusuri bersama.}
\end{frame}

\begin{frame}{Error yang sering muncul minggu ini}
\framesubtitle{Pesan diambil langsung dari g++}
\footnotesize
\begin{tabular}{@{}>{\raggedright\arraybackslash}p{208pt}>{\raggedright\arraybackslash}p{256pt}>{\raggedright\arraybackslash}p{398pt}@{}}
\kepala{Kesalahan} & \kepala{Contoh} & \kepala{Pesan pertama dari g++}\\
Lupa ; sesudah struct & \texttt{struct A \{ int x; \}} & \texttt{error: expected ';' after struct definition}\\
\barisgenap Field tidak ada & \texttt{m.umur pada Mahasiswa} & \texttt{error: 'struct Mahasiswa' has no member named 'umur'}\\
Akses field tanpa indeks & \texttt{rak.judul pada array} & \texttt{error: request for member 'judul' in 'rak', which is of non-class type 'Buku ...}\\
\barisgenap Rekursi tanpa return di satu jalur & \texttt{if (n == 0) return 1; tanpa return lain} & \texttt{warning: control reaches end of non-void function}\\
Rekursi tanpa basis & \texttt{return n * f(n - 1); saja} & \texttt{warning: infinite recursion detected}\\
\garisdasar
\end{tabular}
\vspace{10pt}

{\small Pesan di tabel ini dihasilkan dengan g++ dan opsi -Wall, sama seperti di cek.sh.}
\note{Minta mahasiswa memotret slide ini.}
\end{frame}

\Bagian{4}{Hands-on bertingkat}{45 menit, dari dasar sampai tantangan}
\note{Mahasiswa membuka folder 02 Hands-on/P14 - Rekursi dan Struct - Hands-on.}

\begin{frame}{Peta hands-on}
\framesubtitle{Folder 02 Hands-on/P14 - Rekursi dan Struct - Hands-on}
\small
\begin{tabular}{@{}>{\raggedright\arraybackslash}p{36pt}>{\raggedright\arraybackslash}p{267pt}>{\raggedright\arraybackslash}p{110pt}>{\raggedright\arraybackslash}p{83pt}>{\raggedright\arraybackslash}p{341pt}@{}}
\kepala{No} & \kepala{File} & \kepala{Tingkat} & \kepala{Waktu} & \kepala{Yang dilatih}\\
1 & \texttt{handson1-rekursi.cpp} & Dasar & 10' & kasus basis dan langkah rekursif\\
\barisgenap 2 & \texttt{handson2-struct.cpp} & Menengah & 15' & \texttt{struct}, array of struct, maksimum\\
3 & \texttt{handson3-basis.cpp} & Tantangan & 20' & pohon pemanggilan, perbaiki tiga kesalahan\\
\barisgenap + & \texttt{bonus-hanoi.cpp} & Pengayaan & sisa & rekursi bercabang klasik\\
\garisdasar
\end{tabular}
\vspace{10pt}

Kerjakan berurutan. Cocokkan keluaran dengan folder \texttt{expected/}; latihan yang membaca masukan memakai data di folder \texttt{input/}.\note{Saat berkeliling, minta mahasiswa yang mengerjakan latihan tantangan menjelaskan blok PENJELASAN mereka.}
\end{frame}

\begin{frame}{Cara mengecek dan aturannya}
\framesubtitle{Hands-on}
\begin{columns}[T]
\begin{column}{0.47\textwidth}
\begin{Blok}{Di GDB Online}
Salin isi file latihan, lengkapi TODO, klik Run. Bila program membaca masukan, ketik isi file di \texttt{input/}. Bandingkan dengan \texttt{expected/}.
\end{Blok}
\end{column}
\begin{column}{0.47\textwidth}
\begin{BlokGelap}{Di komputer sendiri}
Jalankan \texttt{bash cek.sh}. Skrip mengompilasi dengan \texttt{-Wall -Werror}, memberi masukan dari \texttt{input/}, lalu mencocokkan keluaran.
\end{BlokGelap}
\end{column}
\end{columns}
\vspace{6pt}
\begin{Catatan}
AI boleh untuk memahami pesan error, tidak untuk meminta solusi utuh. Exit Ticket dikerjakan tanpa bantuan apa pun.
\end{Catatan}
\note{Hands-on tidak dinilai langsung; yang dinilai Exit Ticket dan tugas.}
\end{frame}

\Bagian{5}{Exit Ticket dan tugas}{25 menit terakhir, lalu pekerjaan rumah}
\note{Mulai Exit Ticket tepat waktu.}

\begin{frame}{Exit Ticket 14}
\framesubtitle{Individual, 25 menit, tanpa AI dan internet}
\small
\begin{tabular}{@{}>{\raggedright\arraybackslash}p{60pt}>{\raggedright\arraybackslash}p{560pt}>{\raggedright\arraybackslash}p{80pt}>{\raggedright\arraybackslash}p{150pt}@{}}
\kepala{Soal} & \kepala{Isi} & \kepala{Poin} & \kepala{CPMK}\\
A1 & Tulis fungsi rekursif yang menghitung jumlah 1 sampai n & 25 & CPMK0607\\
\barisgenap A2 & Definisikan struct dan tampilkan data dengan nilai field tertinggi dari array of struct & 25 & CPMK0607\\
B1 & Gambar pohon pemanggilan sebuah fungsi rekursif dan tentukan hasilnya & 20 & CPMK0608\\
\barisgenap B2 & Jelaskan apa yang terjadi bila kasus basis sebuah fungsi dihapus & 20 & CPMK0608\\
B3 & Jelaskan keunggulan array of struct dibandingkan beberapa array sejajar & 10 & CPMK0608\\
\garisdasar
\end{tabular}
\vspace{10pt}

{\small Soal A dinilai dari keluaran yang tepat sama. Soal B dinilai dari ketepatan dan alasan. Pelanggaran aturan membuat nilai Exit Ticket ini 0.}
\note{Soal lengkap dibagikan terpisah dan tidak disimpan di repo publik.}
\end{frame}

\begin{frame}{Tugas 14: Perpustakaan Mini dengan Struct dan Rekursi}
\framesubtitle{Dikumpulkan sebelum Pertemuan 15}
\begin{columns}[T]
\begin{column}{0.47\textwidth}
\begin{Blok}{Bagian A, program, 50 poin}
Buat program perpustakaan mini dengan ketentuan berikut. Semua fungsi pencarian dan agregasi wajib rekursif, bukan perulangan. Ketentuannya di slide berikut.
\end{Blok}
\end{column}
\begin{column}{0.47\textwidth}
\begin{BlokGelap}{Bagian B, penjelasan, 50 poin}
Komentar maksimal 150 kata: gambar tumpukan pemanggilan \texttt{hitungTotalStok} untuk tiga buku contoh, jelaskan kasus basis setiap fungsi rekursif, dan alasan memakai struct dibandingkan array sejajar.
\end{BlokGelap}
\end{column}
\end{columns}
\vspace{8pt}
{\small Data di tugas adalah data kalian sendiri. Rubrik lengkap ada di Modul Belajar 14.}
\note{Ingatkan bahwa penjelasan disalin dari AI mudah dikenali karena tidak menyebut pengalaman mereka sendiri.}
\end{frame}

\begin{frame}{Ketentuan Tugas 14}
\framesubtitle{Perpustakaan Mini dengan Struct dan Rekursi}
\footnotesize
\begin{tabular}{@{}>{\raggedright\arraybackslash}p{87pt}>{\raggedright\arraybackslash}p{788pt}@{}}
\kepala{No} & \kepala{Ketentuan Bagian A}\\
1 & \texttt{struct Buku} berisi kodeBuku, judul, penulis, tahun, dan stok\\
\barisgenap 2 & Array of struct dengan maksimal 10 buku dan validasi masukan\\
3 & Fungsi rekursif \texttt{hitungTotalStok(Buku arr[], int n)}\\
\barisgenap 4 & Fungsi rekursif \texttt{cariBukuTerlama(Buku arr[], int n)} yang mengembalikan indeks buku dengan tahun terkecil\\
5 & Fungsi rekursif \texttt{faktorial(int n)} untuk menghitung banyak cara menyusun n buku di rak\\
\barisgenap 6 & Menu: tambah buku, tampilkan semua, cari buku terlama, total stok, banyak susunan, keluar\\
Bonus & Tambahkan pencarian judul secara rekursif.\\
\garisdasar
\end{tabular}
\note{Bacakan ketentuan satu per satu dan tanyakan apakah ada yang belum jelas.}
\end{frame}

\begin{frame}{Rubrik Tugas 14}
\framesubtitle{Empat level, sama di semua instrumen}
\footnotesize
\begin{tabular}{@{}>{\raggedright\arraybackslash}p{166pt}>{\raggedright\arraybackslash}p{44pt}>{\raggedright\arraybackslash}p{166pt}>{\raggedright\arraybackslash}p{156pt}>{\raggedright\arraybackslash}p{146pt}>{\raggedright\arraybackslash}p{146pt}@{}}
\kepala{Kriteria} & \kepala{Poin} & \kepala{Sangat baik 85--100} & \kepala{Baik 70--84} & \kepala{Cukup 55--69} & \kepala{Kurang <55}\\
A1 Program berjalan & 20 & tanpa error dan peringatan & ada peringatan & perlu satu perbaikan & tidak terkompilasi\\
\barisgenap A2 Struct dan fungsi rekursif & 15 & struct lengkap, tiga fungsi rekursif & satu fungsi tidak rekursif & dua tidak rekursif & tanpa struct\\
A3 Menu dan ketepatan & 15 & semua menu tepat, validasi ada & satu salah & dua salah & sebagian besar salah\\
\barisgenap B1 Tumpukan pemanggilan dan basis & 30 & jejak tepat, basis dijelaskan & satu langkah keliru & tanpa jejak & tidak ada atau disalin\\
B2 Cerita error dan pengujian & 20 & pesan, sebab, perbaikan, uji & pesan tidak dikutip & hanya perbaikan & tidak ada\\
\garisdasar
\end{tabular}
\vspace{8pt}

{\small Nilai kriteria = poin × skor level. Bagian A total 50, Bagian B total 50.}
\note{Rubrik ini sama strukturnya setiap minggu; yang berubah hanya isi kriteria A2, A3, dan B1.}
\end{frame}

\begin{frame}{Sebelum Pertemuan 15}
\framesubtitle{Persiapan}
\begin{columns}[T]
\begin{column}{0.31\textwidth}
\begin{Langkah}{1}{Selesaikan}
Hands-on yang belum lulus \texttt{cek.sh}, terutama latihan tantangan.
\end{Langkah}
\end{column}
\begin{column}{0.31\textwidth}
\begin{Langkah}{2}{Kumpulkan}
Tugas 14 sebelum Pertemuan 15 dimulai.
\end{Langkah}
\end{column}
\begin{column}{0.31\textwidth}
\begin{Langkah}{3}{Baca}
Modul Belajar 14 untuk mengulang, lalu bagian awal modul berikutnya.
\end{Langkah}
\end{column}
\end{columns}
\vspace{10pt}
Pertemuan 15 adalah review dan simulasi UAS untuk materi Pertemuan 9 sampai 14. Bawa catatan dan daftar hal yang masih membingungkan.\note{Tanyakan satu hal yang paling membingungkan hari ini untuk dibuka di review minggu depan.}
\end{frame}

\SlidePenutup{Terima kasih}{Sampai jumpa di Pertemuan 15: review dan simulasi UAS}
\note{Slide penutup.}
\end{document}
"
\documentclass[t]{beamer}
\usetheme{ITERA}
% Mode catatan pembicara: aktif bila \catatan didefinisikan sebelum \input.
\ifdefined\catatan
  \setbeameroption{show only notes}
  \setbeamertemplate{note page}{\vspace*{20pt}\hspace*{4pt}\begin{minipage}{900pt}
    {\ITERAKickerFont\fontsize{11pt}{14pt}\selectfont\color{ITERAGoldText}CATATAN PEMBICARA \ \ |\ \ SLIDE \insertframenumber}\par\vspace{4pt}
    {\ITERADisplay\bfseries\fontsize{22pt}{28pt}\selectfont\color{ITERAStructure}\insertframetitle}\par\vspace{12pt}
    \fontsize{15pt}{23pt}\selectfont\insertnote\end{minipage}}
\fi
\tikzset{fase/.style={draw=none,fill=#1,rounded corners=3pt,minimum width=158pt,minimum height=118pt,text width=146pt,align=center}}

\renewcommand\ITERAmeeting{Pertemuan 14}
\begin{document}
\SlideJudul{IF25-11002 PRAKTIKUM PEMROGRAMAN}{Pertemuan 14\\Rekursi dan Struct}{Fungsi yang memanggil dirinya sendiri, dan tipe data buatan sendiri untuk mengelompokkan data}{Tim Pengajar}{Tahun 2026. Sejajar dengan Pertemuan 14 Algoritma Pemrograman: Rekursi dan Struct.}
\note{Buka dengan menanyakan satu hal dari pertemuan lalu yang masih membingungkan.}

\begin{frame}{Dua mata kuliah, satu minggu}
\framesubtitle{Sebelum mulai}
Topik minggu ini sama di dua mata kuliah, dengan peran yang berbeda.
\vspace{10pt}

\begin{columns}[T]
\begin{column}{0.47\textwidth}
\begin{Blok}{Algoritma: Rekursi dan Struct}
Definisi rekursif, kasus basis, pohon rekursi, serta record atau tipe bentukan sebagai kumpulan field.
\end{Blok}
\end{column}
\begin{column}{0.47\textwidth}
\begin{BlokGelap}{Praktikum: Rekursi dan struct}
Menulis fungsi rekursif dengan kasus basis yang tepat, mendefinisikan \texttt{struct}, dan memproses array of struct.
\end{BlokGelap}
\end{column}
\end{columns}
\note{Tanyakan apa yang sudah dibahas di Algoritma minggu ini. Gunakan jawabannya sebagai jembatan ke praktikum.}
\end{frame}

\begin{frame}{Saat pulang nanti, kalian bisa}
\framesubtitle{Target hari ini}
\begin{columns}[T]
\begin{column}{0.47\textwidth}
\begin{Langkah}{1}{Menerapkan}
Menulis fungsi rekursif dengan kasus basis dan langkah rekursif yang tepat, serta mendefinisikan \texttt{struct} dan memproses array of struct

{\footnotesize\color{ITERAInkMuted}Sub-CPMK 14.1, CPMK0607}
\end{Langkah}
\end{column}
\begin{column}{0.47\textwidth}
\begin{Langkah}{2}{Menjelaskan}
Menelusuri pemanggilan rekursif dengan pohon pemanggilan atau tumpukan panggilan, dan menjelaskan peran kasus basis serta akses anggota struct

{\footnotesize\color{ITERAInkMuted}Sub-CPMK 14.2, CPMK0608}
\end{Langkah}
\end{column}
\end{columns}
\vspace{8pt}
\begin{Catatan}
Dua kemampuan ini dinilai sama berat di Exit Ticket, tugas, UTS, dan UAS.
\end{Catatan}
\note{Bacakan target dengan pelan. Kata kuncinya menerapkan dan menjelaskan.}
\end{frame}

\begin{frame}{Ingat minggu lalu}
\framesubtitle{Review, 5 menit}
\begin{columns}[T]
\begin{column}{0.31\textwidth}
\begin{BlokGelap}{Pertanyaan 1}
Apa beda pass by value dan pass by reference?
\end{BlokGelap}
\end{column}
\begin{column}{0.31\textwidth}
\begin{BlokGelap}{Pertanyaan 2}
Mengapa fungsi non-void harus punya \texttt{return} di semua jalur?
\end{BlokGelap}
\end{column}
\begin{column}{0.31\textwidth}
\begin{BlokGelap}{Pertanyaan 3}
Apakah array di dalam fungsi adalah salinan?
\end{BlokGelap}
\end{column}
\end{columns}
\vspace{8pt}
Jawab di kertas dulu, lalu cocokkan dengan teman sebelah.\note{Pilih dua mahasiswa untuk menjawab. Koreksi singkat, jangan mengajar ulang.}
\end{frame}

\Bagian{1}{Masalah pembuka}{Data buku perpustakaan}
\note{Masalah pembuka 10 menit.}

\begin{frame}[fragile]{Data buku perpustakaan}
\framesubtitle{Masalah minggu ini}
\begin{columns}[T]
\begin{column}{0.55\textwidth}
{\small Perpustakaan prodi mencatat setiap buku dengan kode, judul, penulis, tahun, dan stok. Dengan materi sejauh ini, kalian butuh lima array sejajar dan harus menjaga indeksnya selalu sama.

\vspace{6pt}
Selain itu, beberapa pertanyaan lebih mudah dirumuskan dari ukuran yang lebih kecil, misalnya: total stok n buku adalah stok buku terakhir ditambah total stok n − 1 buku sebelumnya.\par}
\vspace{6pt}
\begin{Catatan}
\textbf{Diskusi 2 menit.} Bagaimana menyimpan lima keterangan satu buku dalam satu variabel? Bagaimana menghitung total stok tanpa perulangan?
\end{Catatan}
\end{column}
\begin{column}{0.41\textwidth}
\begin{Keluaran}[]
[B01] Algoritma (2019) stok 5
[B02] C++ Dasar (2021) stok 8
[B03] Struktur Data (2018) stok 10
Total stok: 23
Buku terlama: Struktur Data
\end{Keluaran}
\end{column}
\end{columns}
\note{Tampung beberapa jawaban tanpa menilai benar salah.}
\end{frame}

\begin{frame}{Dari langkah ke instruksi}
\framesubtitle{Sambungan dengan Algoritma}
\begin{columns}[T]
\begin{column}{0.31\textwidth}
\begin{Langkah}{1}{Kelompokkan}
Satu buku punya beberapa field: kode, judul, penulis, tahun, stok.
\end{Langkah}
\end{column}
\begin{column}{0.31\textwidth}
\begin{Langkah}{2}{Perkecil masalah}
Total n buku = buku terakhir + total n − 1 buku; total 0 buku = 0.
\end{Langkah}
\end{column}
\begin{column}{0.31\textwidth}
\begin{Langkah}{3}{Terjemahkan}
\texttt{struct Buku} dan fungsi rekursif dengan kasus basis.
\end{Langkah}
\end{column}
\end{columns}
\vspace{10pt}
Langkah 1 dan 2 dilatih di Algoritma. Langkah 3 kita kerjakan hari ini dengan C++.\note{Hubungkan eksplisit ke Algoritma minggu ini.}
\end{frame}

\Bagian{2}{Coba dulu, baru dijelaskan}{25 menit eksperimen di GDB Online}
\note{Struggle 25 menit. Berkeliling, jangan menjelaskan materi dulu.}

\begin{frame}{Misi kalian}
\framesubtitle{Struggle, 25 menit}
\begin{columns}[T]
\begin{column}{0.31\textwidth}
\begin{Langkah}{1}{Hitung mundur}
Tulis fungsi yang menampilkan n, n − 1, sampai 1 tanpa \texttt{for} dan \texttt{while}.
\end{Langkah}
\end{column}
\begin{column}{0.31\textwidth}
\begin{Langkah}{2}{Berhenti}
Apa yang terjadi bila fungsi itu tidak punya syarat berhenti?
\end{Langkah}
\end{column}
\begin{column}{0.31\textwidth}
\begin{Langkah}{3}{Kelompokkan}
Simpan nama, NIM, dan IPK satu mahasiswa dalam satu variabel.
\end{Langkah}
\end{column}
\end{columns}
\vspace{8pt}
\begin{columns}[T]
\begin{column}{0.47\textwidth}
\begin{Blok}{Boleh}
Bertanya ke teman, mengubah apa saja, sengaja membuat error.
\end{Blok}
\end{column}
\begin{column}{0.47\textwidth}
\begin{BlokAlert}{Belum dulu}
Mencari jawaban jadi di internet atau AI.
\end{BlokAlert}
\end{column}
\end{columns}
\note{Catat kesalahan yang sering muncul untuk dibahas di debrief.}
\end{frame}

\begin{frame}{Apa yang kalian temukan?}
\framesubtitle{Debrief struggle}
\begin{columns}[T]
\begin{column}{0.31\textwidth}
\begin{BlokGelap}{Pertanyaan 1}
Bagaimana fungsi bisa memanggil dirinya sendiri?
\end{BlokGelap}
\end{column}
\begin{column}{0.31\textwidth}
\begin{BlokGelap}{Pertanyaan 2}
Syarat apa yang membuatnya berhenti?
\end{BlokGelap}
\end{column}
\begin{column}{0.31\textwidth}
\begin{BlokGelap}{Pertanyaan 3}
Apa yang kalian coba untuk menyatukan tiga data mahasiswa?
\end{BlokGelap}
\end{column}
\end{columns}
\vspace{10pt}
Jawaban kalian akan kita uji di bagian berikutnya.\note{Tulis jawaban di papan; buktikan atau patahkan di bagian materi.}
\end{frame}

\Bagian{3}{Memanggil diri sendiri dan tipe buatan}{Kasus basis, pohon pemanggilan, struct, dan array of struct}
\note{Materi 40 menit. Tunjukkan setiap contoh langsung di GDB Online.}

\begin{frame}[fragile]{Rekursi: kasus basis dan langkah rekursif}
\framesubtitle{Faktorial}
\begin{columns}[T]
\begin{column}{0.55\textwidth}
\begin{Kode}[]{faktorial.cpp}
int faktorial(int n) {
    if (n == 0) return 1;
    return n * faktorial(n - 1);
}

int main() {
    cout << faktorial(5) << endl;
    cout << faktorial(0) << endl;
    return 0;
}
\end{Kode}
\end{column}
\begin{column}{0.41\textwidth}
\only<1>{\begin{TebakDulu}
{\ITERAKickerFont\fontsize{10pt}{12pt}\selectfont\color{ITERAAlert}TEBAK DULU}\par\vspace{4pt}
Sebelum dijalankan, tulis keluarannya di kertas.
\end{TebakDulu}
}\only<2>{\KeluaranFile[]{keluaran/P14/k001.txt}
\begin{Catatan}
Kasus basis menghentikan rekursi. Langkah rekursif harus selalu mendekati kasus basis.
\end{Catatan}
}
\end{column}
\end{columns}
\note<1>{Minta mahasiswa menebak keluaran sebelum membuka slide berikutnya. Tanyakan satu atau dua tebakan yang berbeda, lalu buktikan.}
\end{frame}

\begin{frame}{Menelusuri faktorial(4)}
\framesubtitle{Tumpukan pemanggilan}
\small
\begin{tabular}{@{}>{\raggedright\arraybackslash}p{249pt}>{\raggedright\arraybackslash}p{307pt}>{\raggedright\arraybackslash}p{307pt}@{}}
\kepala{Pemanggilan} & \kepala{Menunggu} & \kepala{Nilai kembali}\\
faktorial(4) & 4 × faktorial(3) & 4 × 6 = 24\\
\barisgenap faktorial(3) & 3 × faktorial(2) & 3 × 2 = 6\\
faktorial(2) & 2 × faktorial(1) & 2 × 1 = 2\\
\barisgenap faktorial(1) & 1 × faktorial(0) & 1 × 1 = 1\\
faktorial(0) & kasus basis & 1\\
\garisdasar
\end{tabular}
\vspace{10pt}

Pemanggilan turun sampai kasus basis, lalu nilai kembali naik dari bawah ke atas.
\end{frame}

\begin{frame}[fragile]{Rekursi bercabang dan rekursi pada array}
\framesubtitle{Fibonacci dan jumlah array}
\begin{columns}[T]
\begin{column}{0.60\textwidth}
\begin{Kode}[listing options={style=ITERACodeMini}]{fibarray.cpp}
int fib(int n) {
    if (n <= 1) return n;
    return fib(n - 1) + fib(n - 2);
}
int jumlah(int a[], int n) {
    if (n == 0) return 0;
    return a[n - 1] + jumlah(a, n - 1);
}
int main() {
    int a[4] = {5, 8, 10, 2};
    cout << fib(7) << " " << jumlah(a, 4) << endl;
    return 0;
}
\end{Kode}
\end{column}
\begin{column}{0.36\textwidth}
\only<1>{\begin{TebakDulu}
{\ITERAKickerFont\fontsize{10pt}{12pt}\selectfont\color{ITERAAlert}TEBAK DULU}\par\vspace{4pt}
Sebelum dijalankan, tulis keluarannya di kertas.
\end{TebakDulu}
}\only<2>{\KeluaranFile[listing options={style=ITERAPlainKecil}]{keluaran/P14/k002.txt}
\begin{Catatan}
\texttt{fib} memanggil dirinya dua kali, sehingga pohon pemanggilannya melebar. \texttt{jumlah} memproses elemen terakhir lalu sisanya.
\end{Catatan}
}
\end{column}
\end{columns}
\note<1>{Minta mahasiswa menebak keluaran sebelum membuka slide berikutnya. Tanyakan satu atau dua tebakan yang berbeda, lalu buktikan.}
\end{frame}

\begin{frame}[fragile]{struct: tipe data buatan sendiri}
\framesubtitle{Mengelompokkan field}
\begin{columns}[T]
\begin{column}{0.60\textwidth}
\begin{Kode}[listing options={style=ITERACodeMini}]{struct.cpp}
struct Mahasiswa {
    string nama;
    string nim;
    double ipk;
};
int main() {
    Mahasiswa m;
    m.nama = "Rani";
    m.nim = "126140001";
    m.ipk = 3.72;
    cout << m.nama << " " << m.nim << " " << m.ipk << endl;
    return 0;
}
\end{Kode}
\end{column}
\begin{column}{0.36\textwidth}
\only<1>{\begin{TebakDulu}
{\ITERAKickerFont\fontsize{10pt}{12pt}\selectfont\color{ITERAAlert}TEBAK DULU}\par\vspace{4pt}
Sebelum dijalankan, tulis keluarannya di kertas.
\end{TebakDulu}
}\only<2>{\KeluaranFile[listing options={style=ITERAPlainKecil}]{keluaran/P14/k003.txt}
\begin{Catatan}
Definisi \texttt{struct} diakhiri titik koma. Field diakses dengan titik: \texttt{m.nama}.
\end{Catatan}
}
\end{column}
\end{columns}
\note<1>{Minta mahasiswa menebak keluaran sebelum membuka slide berikutnya. Tanyakan satu atau dua tebakan yang berbeda, lalu buktikan.}
\end{frame}

\begin{frame}[fragile]{Array of struct}
\framesubtitle{Banyak data terstruktur}
\begin{columns}[T]
\begin{column}{0.60\textwidth}
\begin{Kode}[listing options={style=ITERACodeKecil}]{arraystruct.cpp}
struct Buku { string judul; int tahun; };

int main() {
    Buku rak[3] = {{"Algoritma", 2019},
                   {"C++ Dasar", 2021},
                   {"Struktur Data", 2018}};
    int iLama = 0;
    for (int i = 1; i < 3; i++)
        if (rak[i].tahun < rak[iLama].tahun) iLama = i;
    cout << rak[iLama].judul << endl;
    return 0;
}
\end{Kode}
\end{column}
\begin{column}{0.36\textwidth}
\only<1>{\begin{TebakDulu}
{\ITERAKickerFont\fontsize{10pt}{12pt}\selectfont\color{ITERAAlert}TEBAK DULU}\par\vspace{4pt}
Sebelum dijalankan, tulis keluarannya di kertas.
\end{TebakDulu}
}\only<2>{\KeluaranFile[listing options={style=ITERAPlainKecil}]{keluaran/P14/k004.txt}
\begin{Catatan}
\texttt{rak[i].tahun}: indeks dulu untuk memilih buku, lalu titik untuk memilih field.
\end{Catatan}
}
\end{column}
\end{columns}
\note<1>{Minta mahasiswa menebak keluaran sebelum membuka slide berikutnya. Tanyakan satu atau dua tebakan yang berbeda, lalu buktikan.}
\end{frame}

\begin{frame}[fragile]{Tebak keluarannya}
\framesubtitle{Latihan menjelaskan, 3 menit}
\begin{columns}[T]
\begin{column}{0.56\textwidth}
\begin{Kode}[]{tebak.cpp}
int f(int n) {
    if (n <= 0) return 0;
    if (n % 2 == 0) return n + f(n - 1);
    return f(n - 1);
}

int main() {
    cout << f(6) << " " << f(3) << endl;
    return 0;
}
\end{Kode}
\end{column}
\begin{column}{0.40\textwidth}
\begin{Catatan}
Tulis keluarannya di kertas, karakter demi karakter. Jangan dijalankan dulu.
\end{Catatan}
\end{column}
\end{columns}
\note{Contoh soal Bagian B. Beri waktu 3 menit sebelum slide jawaban.}
\end{frame}

\begin{frame}[fragile]{Tebak keluarannya: jawaban}
\framesubtitle{Latihan menjelaskan}
\begin{columns}[T]
\begin{column}{0.56\textwidth}
\begin{Kode}[]{tebak.cpp}
int f(int n) {
    if (n <= 0) return 0;
    if (n % 2 == 0) return n + f(n - 1);
    return f(n - 1);
}

int main() {
    cout << f(6) << " " << f(3) << endl;
    return 0;
}
\end{Kode}
\end{column}
\begin{column}{0.40\textwidth}
\begin{Keluaran}[]
12 2
\end{Keluaran}
\begin{Catatan}
Hanya bilangan genap yang dijumlahkan: f(6) = 6 + 4 + 2 = 12, f(3) = 2.
\end{Catatan}
\end{column}
\end{columns}
\note{Tanyakan jawaban yang paling sering salah, lalu telusuri bersama.}
\end{frame}

\begin{frame}{Error yang sering muncul minggu ini}
\framesubtitle{Pesan diambil langsung dari g++}
\footnotesize
\begin{tabular}{@{}>{\raggedright\arraybackslash}p{208pt}>{\raggedright\arraybackslash}p{256pt}>{\raggedright\arraybackslash}p{398pt}@{}}
\kepala{Kesalahan} & \kepala{Contoh} & \kepala{Pesan pertama dari g++}\\
Lupa ; sesudah struct & \texttt{struct A \{ int x; \}} & \texttt{error: expected ';' after struct definition}\\
\barisgenap Field tidak ada & \texttt{m.umur pada Mahasiswa} & \texttt{error: 'struct Mahasiswa' has no member named 'umur'}\\
Akses field tanpa indeks & \texttt{rak.judul pada array} & \texttt{error: request for member 'judul' in 'rak', which is of non-class type 'Buku ...}\\
\barisgenap Rekursi tanpa return di satu jalur & \texttt{if (n == 0) return 1; tanpa return lain} & \texttt{warning: control reaches end of non-void function}\\
Rekursi tanpa basis & \texttt{return n * f(n - 1); saja} & \texttt{warning: infinite recursion detected}\\
\garisdasar
\end{tabular}
\vspace{10pt}

{\small Pesan di tabel ini dihasilkan dengan g++ dan opsi -Wall, sama seperti di cek.sh.}
\note{Minta mahasiswa memotret slide ini.}
\end{frame}

\Bagian{4}{Hands-on bertingkat}{45 menit, dari dasar sampai tantangan}
\note{Mahasiswa membuka folder 02 Hands-on/P14 - Rekursi dan Struct - Hands-on.}

\begin{frame}{Peta hands-on}
\framesubtitle{Folder 02 Hands-on/P14 - Rekursi dan Struct - Hands-on}
\small
\begin{tabular}{@{}>{\raggedright\arraybackslash}p{36pt}>{\raggedright\arraybackslash}p{267pt}>{\raggedright\arraybackslash}p{110pt}>{\raggedright\arraybackslash}p{83pt}>{\raggedright\arraybackslash}p{341pt}@{}}
\kepala{No} & \kepala{File} & \kepala{Tingkat} & \kepala{Waktu} & \kepala{Yang dilatih}\\
1 & \texttt{handson1-rekursi.cpp} & Dasar & 10' & kasus basis dan langkah rekursif\\
\barisgenap 2 & \texttt{handson2-struct.cpp} & Menengah & 15' & \texttt{struct}, array of struct, maksimum\\
3 & \texttt{handson3-basis.cpp} & Tantangan & 20' & pohon pemanggilan, perbaiki tiga kesalahan\\
\barisgenap + & \texttt{bonus-hanoi.cpp} & Pengayaan & sisa & rekursi bercabang klasik\\
\garisdasar
\end{tabular}
\vspace{10pt}

Kerjakan berurutan. Cocokkan keluaran dengan folder \texttt{expected/}; latihan yang membaca masukan memakai data di folder \texttt{input/}.\note{Saat berkeliling, minta mahasiswa yang mengerjakan latihan tantangan menjelaskan blok PENJELASAN mereka.}
\end{frame}

\begin{frame}{Cara mengecek dan aturannya}
\framesubtitle{Hands-on}
\begin{columns}[T]
\begin{column}{0.47\textwidth}
\begin{Blok}{Di GDB Online}
Salin isi file latihan, lengkapi TODO, klik Run. Bila program membaca masukan, ketik isi file di \texttt{input/}. Bandingkan dengan \texttt{expected/}.
\end{Blok}
\end{column}
\begin{column}{0.47\textwidth}
\begin{BlokGelap}{Di komputer sendiri}
Jalankan \texttt{bash cek.sh}. Skrip mengompilasi dengan \texttt{-Wall -Werror}, memberi masukan dari \texttt{input/}, lalu mencocokkan keluaran.
\end{BlokGelap}
\end{column}
\end{columns}
\vspace{6pt}
\begin{Catatan}
AI boleh untuk memahami pesan error, tidak untuk meminta solusi utuh. Exit Ticket dikerjakan tanpa bantuan apa pun.
\end{Catatan}
\note{Hands-on tidak dinilai langsung; yang dinilai Exit Ticket dan tugas.}
\end{frame}

\Bagian{5}{Exit Ticket dan tugas}{25 menit terakhir, lalu pekerjaan rumah}
\note{Mulai Exit Ticket tepat waktu.}

\begin{frame}{Exit Ticket 14}
\framesubtitle{Individual, 25 menit, tanpa AI dan internet}
\small
\begin{tabular}{@{}>{\raggedright\arraybackslash}p{60pt}>{\raggedright\arraybackslash}p{560pt}>{\raggedright\arraybackslash}p{80pt}>{\raggedright\arraybackslash}p{150pt}@{}}
\kepala{Soal} & \kepala{Isi} & \kepala{Poin} & \kepala{CPMK}\\
A1 & Tulis fungsi rekursif yang menghitung jumlah 1 sampai n & 25 & CPMK0607\\
\barisgenap A2 & Definisikan struct dan tampilkan data dengan nilai field tertinggi dari array of struct & 25 & CPMK0607\\
B1 & Gambar pohon pemanggilan sebuah fungsi rekursif dan tentukan hasilnya & 20 & CPMK0608\\
\barisgenap B2 & Jelaskan apa yang terjadi bila kasus basis sebuah fungsi dihapus & 20 & CPMK0608\\
B3 & Jelaskan keunggulan array of struct dibandingkan beberapa array sejajar & 10 & CPMK0608\\
\garisdasar
\end{tabular}
\vspace{10pt}

{\small Soal A dinilai dari keluaran yang tepat sama. Soal B dinilai dari ketepatan dan alasan. Pelanggaran aturan membuat nilai Exit Ticket ini 0.}
\note{Soal lengkap dibagikan terpisah dan tidak disimpan di repo publik.}
\end{frame}

\begin{frame}{Tugas 14: Perpustakaan Mini dengan Struct dan Rekursi}
\framesubtitle{Dikumpulkan sebelum Pertemuan 15}
\begin{columns}[T]
\begin{column}{0.47\textwidth}
\begin{Blok}{Bagian A, program, 50 poin}
Buat program perpustakaan mini dengan ketentuan berikut. Semua fungsi pencarian dan agregasi wajib rekursif, bukan perulangan. Ketentuannya di slide berikut.
\end{Blok}
\end{column}
\begin{column}{0.47\textwidth}
\begin{BlokGelap}{Bagian B, penjelasan, 50 poin}
Komentar maksimal 150 kata: gambar tumpukan pemanggilan \texttt{hitungTotalStok} untuk tiga buku contoh, jelaskan kasus basis setiap fungsi rekursif, dan alasan memakai struct dibandingkan array sejajar.
\end{BlokGelap}
\end{column}
\end{columns}
\vspace{8pt}
{\small Data di tugas adalah data kalian sendiri. Rubrik lengkap ada di Modul Belajar 14.}
\note{Ingatkan bahwa penjelasan disalin dari AI mudah dikenali karena tidak menyebut pengalaman mereka sendiri.}
\end{frame}

\begin{frame}{Ketentuan Tugas 14}
\framesubtitle{Perpustakaan Mini dengan Struct dan Rekursi}
\footnotesize
\begin{tabular}{@{}>{\raggedright\arraybackslash}p{87pt}>{\raggedright\arraybackslash}p{788pt}@{}}
\kepala{No} & \kepala{Ketentuan Bagian A}\\
1 & \texttt{struct Buku} berisi kodeBuku, judul, penulis, tahun, dan stok\\
\barisgenap 2 & Array of struct dengan maksimal 10 buku dan validasi masukan\\
3 & Fungsi rekursif \texttt{hitungTotalStok(Buku arr[], int n)}\\
\barisgenap 4 & Fungsi rekursif \texttt{cariBukuTerlama(Buku arr[], int n)} yang mengembalikan indeks buku dengan tahun terkecil\\
5 & Fungsi rekursif \texttt{faktorial(int n)} untuk menghitung banyak cara menyusun n buku di rak\\
\barisgenap 6 & Menu: tambah buku, tampilkan semua, cari buku terlama, total stok, banyak susunan, keluar\\
Bonus & Tambahkan pencarian judul secara rekursif.\\
\garisdasar
\end{tabular}
\note{Bacakan ketentuan satu per satu dan tanyakan apakah ada yang belum jelas.}
\end{frame}

\begin{frame}{Rubrik Tugas 14}
\framesubtitle{Empat level, sama di semua instrumen}
\footnotesize
\begin{tabular}{@{}>{\raggedright\arraybackslash}p{166pt}>{\raggedright\arraybackslash}p{44pt}>{\raggedright\arraybackslash}p{166pt}>{\raggedright\arraybackslash}p{156pt}>{\raggedright\arraybackslash}p{146pt}>{\raggedright\arraybackslash}p{146pt}@{}}
\kepala{Kriteria} & \kepala{Poin} & \kepala{Sangat baik 85--100} & \kepala{Baik 70--84} & \kepala{Cukup 55--69} & \kepala{Kurang <55}\\
A1 Program berjalan & 20 & tanpa error dan peringatan & ada peringatan & perlu satu perbaikan & tidak terkompilasi\\
\barisgenap A2 Struct dan fungsi rekursif & 15 & struct lengkap, tiga fungsi rekursif & satu fungsi tidak rekursif & dua tidak rekursif & tanpa struct\\
A3 Menu dan ketepatan & 15 & semua menu tepat, validasi ada & satu salah & dua salah & sebagian besar salah\\
\barisgenap B1 Tumpukan pemanggilan dan basis & 30 & jejak tepat, basis dijelaskan & satu langkah keliru & tanpa jejak & tidak ada atau disalin\\
B2 Cerita error dan pengujian & 20 & pesan, sebab, perbaikan, uji & pesan tidak dikutip & hanya perbaikan & tidak ada\\
\garisdasar
\end{tabular}
\vspace{8pt}

{\small Nilai kriteria = poin × skor level. Bagian A total 50, Bagian B total 50.}
\note{Rubrik ini sama strukturnya setiap minggu; yang berubah hanya isi kriteria A2, A3, dan B1.}
\end{frame}

\begin{frame}{Sebelum Pertemuan 15}
\framesubtitle{Persiapan}
\begin{columns}[T]
\begin{column}{0.31\textwidth}
\begin{Langkah}{1}{Selesaikan}
Hands-on yang belum lulus \texttt{cek.sh}, terutama latihan tantangan.
\end{Langkah}
\end{column}
\begin{column}{0.31\textwidth}
\begin{Langkah}{2}{Kumpulkan}
Tugas 14 sebelum Pertemuan 15 dimulai.
\end{Langkah}
\end{column}
\begin{column}{0.31\textwidth}
\begin{Langkah}{3}{Baca}
Modul Belajar 14 untuk mengulang, lalu bagian awal modul berikutnya.
\end{Langkah}
\end{column}
\end{columns}
\vspace{10pt}
Pertemuan 15 adalah review dan simulasi UAS untuk materi Pertemuan 9 sampai 14. Bawa catatan dan daftar hal yang masih membingungkan.\note{Tanyakan satu hal yang paling membingungkan hari ini untuk dibuka di review minggu depan.}
\end{frame}

\SlidePenutup{Terima kasih}{Sampai jumpa di Pertemuan 15: review dan simulasi UAS}
\note{Slide penutup.}
\end{document}
"
\documentclass[t]{beamer}
\usetheme{ITERA}
% Mode catatan pembicara: aktif bila \catatan didefinisikan sebelum \input.
\ifdefined\catatan
  \setbeameroption{show only notes}
  \setbeamertemplate{note page}{\vspace*{20pt}\hspace*{4pt}\begin{minipage}{900pt}
    {\ITERAKickerFont\fontsize{11pt}{14pt}\selectfont\color{ITERAGoldText}CATATAN PEMBICARA \ \ |\ \ SLIDE \insertframenumber}\par\vspace{4pt}
    {\ITERADisplay\bfseries\fontsize{22pt}{28pt}\selectfont\color{ITERAStructure}\insertframetitle}\par\vspace{12pt}
    \fontsize{15pt}{23pt}\selectfont\insertnote\end{minipage}}
\fi
\tikzset{fase/.style={draw=none,fill=#1,rounded corners=3pt,minimum width=158pt,minimum height=118pt,text width=146pt,align=center}}

\renewcommand\ITERAmeeting{Pertemuan 14}
\begin{document}
\SlideJudul{IF25-11002 PRAKTIKUM PEMROGRAMAN}{Pertemuan 14\\Rekursi dan Struct}{Fungsi yang memanggil dirinya sendiri, dan tipe data buatan sendiri untuk mengelompokkan data}{Tim Pengajar}{Tahun 2026. Sejajar dengan Pertemuan 14 Algoritma Pemrograman: Rekursi dan Struct.}
\note{Buka dengan menanyakan satu hal dari pertemuan lalu yang masih membingungkan.}

\begin{frame}{Dua mata kuliah, satu minggu}
\framesubtitle{Sebelum mulai}
Topik minggu ini sama di dua mata kuliah, dengan peran yang berbeda.
\vspace{10pt}

\begin{columns}[T]
\begin{column}{0.47\textwidth}
\begin{Blok}{Algoritma: Rekursi dan Struct}
Definisi rekursif, kasus basis, pohon rekursi, serta record atau tipe bentukan sebagai kumpulan field.
\end{Blok}
\end{column}
\begin{column}{0.47\textwidth}
\begin{BlokGelap}{Praktikum: Rekursi dan struct}
Menulis fungsi rekursif dengan kasus basis yang tepat, mendefinisikan \texttt{struct}, dan memproses array of struct.
\end{BlokGelap}
\end{column}
\end{columns}
\note{Tanyakan apa yang sudah dibahas di Algoritma minggu ini. Gunakan jawabannya sebagai jembatan ke praktikum.}
\end{frame}

\begin{frame}{Saat pulang nanti, kalian bisa}
\framesubtitle{Target hari ini}
\begin{columns}[T]
\begin{column}{0.47\textwidth}
\begin{Langkah}{1}{Menerapkan}
Menulis fungsi rekursif dengan kasus basis dan langkah rekursif yang tepat, serta mendefinisikan \texttt{struct} dan memproses array of struct

{\footnotesize\color{ITERAInkMuted}Sub-CPMK 14.1, CPMK0607}
\end{Langkah}
\end{column}
\begin{column}{0.47\textwidth}
\begin{Langkah}{2}{Menjelaskan}
Menelusuri pemanggilan rekursif dengan pohon pemanggilan atau tumpukan panggilan, dan menjelaskan peran kasus basis serta akses anggota struct

{\footnotesize\color{ITERAInkMuted}Sub-CPMK 14.2, CPMK0608}
\end{Langkah}
\end{column}
\end{columns}
\vspace{8pt}
\begin{Catatan}
Dua kemampuan ini dinilai sama berat di Exit Ticket, tugas, UTS, dan UAS.
\end{Catatan}
\note{Bacakan target dengan pelan. Kata kuncinya menerapkan dan menjelaskan.}
\end{frame}

\begin{frame}{Ingat minggu lalu}
\framesubtitle{Review, 5 menit}
\begin{columns}[T]
\begin{column}{0.31\textwidth}
\begin{BlokGelap}{Pertanyaan 1}
Apa beda pass by value dan pass by reference?
\end{BlokGelap}
\end{column}
\begin{column}{0.31\textwidth}
\begin{BlokGelap}{Pertanyaan 2}
Mengapa fungsi non-void harus punya \texttt{return} di semua jalur?
\end{BlokGelap}
\end{column}
\begin{column}{0.31\textwidth}
\begin{BlokGelap}{Pertanyaan 3}
Apakah array di dalam fungsi adalah salinan?
\end{BlokGelap}
\end{column}
\end{columns}
\vspace{8pt}
Jawab di kertas dulu, lalu cocokkan dengan teman sebelah.\note{Pilih dua mahasiswa untuk menjawab. Koreksi singkat, jangan mengajar ulang.}
\end{frame}

\Bagian{1}{Masalah pembuka}{Data buku perpustakaan}
\note{Masalah pembuka 10 menit.}

\begin{frame}[fragile]{Data buku perpustakaan}
\framesubtitle{Masalah minggu ini}
\begin{columns}[T]
\begin{column}{0.55\textwidth}
{\small Perpustakaan prodi mencatat setiap buku dengan kode, judul, penulis, tahun, dan stok. Dengan materi sejauh ini, kalian butuh lima array sejajar dan harus menjaga indeksnya selalu sama.

\vspace{6pt}
Selain itu, beberapa pertanyaan lebih mudah dirumuskan dari ukuran yang lebih kecil, misalnya: total stok n buku adalah stok buku terakhir ditambah total stok n − 1 buku sebelumnya.\par}
\vspace{6pt}
\begin{Catatan}
\textbf{Diskusi 2 menit.} Bagaimana menyimpan lima keterangan satu buku dalam satu variabel? Bagaimana menghitung total stok tanpa perulangan?
\end{Catatan}
\end{column}
\begin{column}{0.41\textwidth}
\begin{Keluaran}[]
[B01] Algoritma (2019) stok 5
[B02] C++ Dasar (2021) stok 8
[B03] Struktur Data (2018) stok 10
Total stok: 23
Buku terlama: Struktur Data
\end{Keluaran}
\end{column}
\end{columns}
\note{Tampung beberapa jawaban tanpa menilai benar salah.}
\end{frame}

\begin{frame}{Dari langkah ke instruksi}
\framesubtitle{Sambungan dengan Algoritma}
\begin{columns}[T]
\begin{column}{0.31\textwidth}
\begin{Langkah}{1}{Kelompokkan}
Satu buku punya beberapa field: kode, judul, penulis, tahun, stok.
\end{Langkah}
\end{column}
\begin{column}{0.31\textwidth}
\begin{Langkah}{2}{Perkecil masalah}
Total n buku = buku terakhir + total n − 1 buku; total 0 buku = 0.
\end{Langkah}
\end{column}
\begin{column}{0.31\textwidth}
\begin{Langkah}{3}{Terjemahkan}
\texttt{struct Buku} dan fungsi rekursif dengan kasus basis.
\end{Langkah}
\end{column}
\end{columns}
\vspace{10pt}
Langkah 1 dan 2 dilatih di Algoritma. Langkah 3 kita kerjakan hari ini dengan C++.\note{Hubungkan eksplisit ke Algoritma minggu ini.}
\end{frame}

\Bagian{2}{Coba dulu, baru dijelaskan}{25 menit eksperimen di GDB Online}
\note{Struggle 25 menit. Berkeliling, jangan menjelaskan materi dulu.}

\begin{frame}{Misi kalian}
\framesubtitle{Struggle, 25 menit}
\begin{columns}[T]
\begin{column}{0.31\textwidth}
\begin{Langkah}{1}{Hitung mundur}
Tulis fungsi yang menampilkan n, n − 1, sampai 1 tanpa \texttt{for} dan \texttt{while}.
\end{Langkah}
\end{column}
\begin{column}{0.31\textwidth}
\begin{Langkah}{2}{Berhenti}
Apa yang terjadi bila fungsi itu tidak punya syarat berhenti?
\end{Langkah}
\end{column}
\begin{column}{0.31\textwidth}
\begin{Langkah}{3}{Kelompokkan}
Simpan nama, NIM, dan IPK satu mahasiswa dalam satu variabel.
\end{Langkah}
\end{column}
\end{columns}
\vspace{8pt}
\begin{columns}[T]
\begin{column}{0.47\textwidth}
\begin{Blok}{Boleh}
Bertanya ke teman, mengubah apa saja, sengaja membuat error.
\end{Blok}
\end{column}
\begin{column}{0.47\textwidth}
\begin{BlokAlert}{Belum dulu}
Mencari jawaban jadi di internet atau AI.
\end{BlokAlert}
\end{column}
\end{columns}
\note{Catat kesalahan yang sering muncul untuk dibahas di debrief.}
\end{frame}

\begin{frame}{Apa yang kalian temukan?}
\framesubtitle{Debrief struggle}
\begin{columns}[T]
\begin{column}{0.31\textwidth}
\begin{BlokGelap}{Pertanyaan 1}
Bagaimana fungsi bisa memanggil dirinya sendiri?
\end{BlokGelap}
\end{column}
\begin{column}{0.31\textwidth}
\begin{BlokGelap}{Pertanyaan 2}
Syarat apa yang membuatnya berhenti?
\end{BlokGelap}
\end{column}
\begin{column}{0.31\textwidth}
\begin{BlokGelap}{Pertanyaan 3}
Apa yang kalian coba untuk menyatukan tiga data mahasiswa?
\end{BlokGelap}
\end{column}
\end{columns}
\vspace{10pt}
Jawaban kalian akan kita uji di bagian berikutnya.\note{Tulis jawaban di papan; buktikan atau patahkan di bagian materi.}
\end{frame}

\Bagian{3}{Memanggil diri sendiri dan tipe buatan}{Kasus basis, pohon pemanggilan, struct, dan array of struct}
\note{Materi 40 menit. Tunjukkan setiap contoh langsung di GDB Online.}

\begin{frame}[fragile]{Rekursi: kasus basis dan langkah rekursif}
\framesubtitle{Faktorial}
\begin{columns}[T]
\begin{column}{0.55\textwidth}
\begin{Kode}[]{faktorial.cpp}
int faktorial(int n) {
    if (n == 0) return 1;
    return n * faktorial(n - 1);
}

int main() {
    cout << faktorial(5) << endl;
    cout << faktorial(0) << endl;
    return 0;
}
\end{Kode}
\end{column}
\begin{column}{0.41\textwidth}
\only<1>{\begin{TebakDulu}
{\ITERAKickerFont\fontsize{10pt}{12pt}\selectfont\color{ITERAAlert}TEBAK DULU}\par\vspace{4pt}
Sebelum dijalankan, tulis keluarannya di kertas.
\end{TebakDulu}
}\only<2>{\KeluaranFile[]{keluaran/P14/k001.txt}
\begin{Catatan}
Kasus basis menghentikan rekursi. Langkah rekursif harus selalu mendekati kasus basis.
\end{Catatan}
}
\end{column}
\end{columns}
\note<1>{Minta mahasiswa menebak keluaran sebelum membuka slide berikutnya. Tanyakan satu atau dua tebakan yang berbeda, lalu buktikan.}
\end{frame}

\begin{frame}{Menelusuri faktorial(4)}
\framesubtitle{Tumpukan pemanggilan}
\small
\begin{tabular}{@{}>{\raggedright\arraybackslash}p{249pt}>{\raggedright\arraybackslash}p{307pt}>{\raggedright\arraybackslash}p{307pt}@{}}
\kepala{Pemanggilan} & \kepala{Menunggu} & \kepala{Nilai kembali}\\
faktorial(4) & 4 × faktorial(3) & 4 × 6 = 24\\
\barisgenap faktorial(3) & 3 × faktorial(2) & 3 × 2 = 6\\
faktorial(2) & 2 × faktorial(1) & 2 × 1 = 2\\
\barisgenap faktorial(1) & 1 × faktorial(0) & 1 × 1 = 1\\
faktorial(0) & kasus basis & 1\\
\garisdasar
\end{tabular}
\vspace{10pt}

Pemanggilan turun sampai kasus basis, lalu nilai kembali naik dari bawah ke atas.
\end{frame}

\begin{frame}[fragile]{Rekursi bercabang dan rekursi pada array}
\framesubtitle{Fibonacci dan jumlah array}
\begin{columns}[T]
\begin{column}{0.60\textwidth}
\begin{Kode}[listing options={style=ITERACodeMini}]{fibarray.cpp}
int fib(int n) {
    if (n <= 1) return n;
    return fib(n - 1) + fib(n - 2);
}
int jumlah(int a[], int n) {
    if (n == 0) return 0;
    return a[n - 1] + jumlah(a, n - 1);
}
int main() {
    int a[4] = {5, 8, 10, 2};
    cout << fib(7) << " " << jumlah(a, 4) << endl;
    return 0;
}
\end{Kode}
\end{column}
\begin{column}{0.36\textwidth}
\only<1>{\begin{TebakDulu}
{\ITERAKickerFont\fontsize{10pt}{12pt}\selectfont\color{ITERAAlert}TEBAK DULU}\par\vspace{4pt}
Sebelum dijalankan, tulis keluarannya di kertas.
\end{TebakDulu}
}\only<2>{\KeluaranFile[listing options={style=ITERAPlainKecil}]{keluaran/P14/k002.txt}
\begin{Catatan}
\texttt{fib} memanggil dirinya dua kali, sehingga pohon pemanggilannya melebar. \texttt{jumlah} memproses elemen terakhir lalu sisanya.
\end{Catatan}
}
\end{column}
\end{columns}
\note<1>{Minta mahasiswa menebak keluaran sebelum membuka slide berikutnya. Tanyakan satu atau dua tebakan yang berbeda, lalu buktikan.}
\end{frame}

\begin{frame}[fragile]{struct: tipe data buatan sendiri}
\framesubtitle{Mengelompokkan field}
\begin{columns}[T]
\begin{column}{0.60\textwidth}
\begin{Kode}[listing options={style=ITERACodeMini}]{struct.cpp}
struct Mahasiswa {
    string nama;
    string nim;
    double ipk;
};
int main() {
    Mahasiswa m;
    m.nama = "Rani";
    m.nim = "126140001";
    m.ipk = 3.72;
    cout << m.nama << " " << m.nim << " " << m.ipk << endl;
    return 0;
}
\end{Kode}
\end{column}
\begin{column}{0.36\textwidth}
\only<1>{\begin{TebakDulu}
{\ITERAKickerFont\fontsize{10pt}{12pt}\selectfont\color{ITERAAlert}TEBAK DULU}\par\vspace{4pt}
Sebelum dijalankan, tulis keluarannya di kertas.
\end{TebakDulu}
}\only<2>{\KeluaranFile[listing options={style=ITERAPlainKecil}]{keluaran/P14/k003.txt}
\begin{Catatan}
Definisi \texttt{struct} diakhiri titik koma. Field diakses dengan titik: \texttt{m.nama}.
\end{Catatan}
}
\end{column}
\end{columns}
\note<1>{Minta mahasiswa menebak keluaran sebelum membuka slide berikutnya. Tanyakan satu atau dua tebakan yang berbeda, lalu buktikan.}
\end{frame}

\begin{frame}[fragile]{Array of struct}
\framesubtitle{Banyak data terstruktur}
\begin{columns}[T]
\begin{column}{0.60\textwidth}
\begin{Kode}[listing options={style=ITERACodeKecil}]{arraystruct.cpp}
struct Buku { string judul; int tahun; };

int main() {
    Buku rak[3] = {{"Algoritma", 2019},
                   {"C++ Dasar", 2021},
                   {"Struktur Data", 2018}};
    int iLama = 0;
    for (int i = 1; i < 3; i++)
        if (rak[i].tahun < rak[iLama].tahun) iLama = i;
    cout << rak[iLama].judul << endl;
    return 0;
}
\end{Kode}
\end{column}
\begin{column}{0.36\textwidth}
\only<1>{\begin{TebakDulu}
{\ITERAKickerFont\fontsize{10pt}{12pt}\selectfont\color{ITERAAlert}TEBAK DULU}\par\vspace{4pt}
Sebelum dijalankan, tulis keluarannya di kertas.
\end{TebakDulu}
}\only<2>{\KeluaranFile[listing options={style=ITERAPlainKecil}]{keluaran/P14/k004.txt}
\begin{Catatan}
\texttt{rak[i].tahun}: indeks dulu untuk memilih buku, lalu titik untuk memilih field.
\end{Catatan}
}
\end{column}
\end{columns}
\note<1>{Minta mahasiswa menebak keluaran sebelum membuka slide berikutnya. Tanyakan satu atau dua tebakan yang berbeda, lalu buktikan.}
\end{frame}

\begin{frame}[fragile]{Tebak keluarannya}
\framesubtitle{Latihan menjelaskan, 3 menit}
\begin{columns}[T]
\begin{column}{0.56\textwidth}
\begin{Kode}[]{tebak.cpp}
int f(int n) {
    if (n <= 0) return 0;
    if (n % 2 == 0) return n + f(n - 1);
    return f(n - 1);
}

int main() {
    cout << f(6) << " " << f(3) << endl;
    return 0;
}
\end{Kode}
\end{column}
\begin{column}{0.40\textwidth}
\begin{Catatan}
Tulis keluarannya di kertas, karakter demi karakter. Jangan dijalankan dulu.
\end{Catatan}
\end{column}
\end{columns}
\note{Contoh soal Bagian B. Beri waktu 3 menit sebelum slide jawaban.}
\end{frame}

\begin{frame}[fragile]{Tebak keluarannya: jawaban}
\framesubtitle{Latihan menjelaskan}
\begin{columns}[T]
\begin{column}{0.56\textwidth}
\begin{Kode}[]{tebak.cpp}
int f(int n) {
    if (n <= 0) return 0;
    if (n % 2 == 0) return n + f(n - 1);
    return f(n - 1);
}

int main() {
    cout << f(6) << " " << f(3) << endl;
    return 0;
}
\end{Kode}
\end{column}
\begin{column}{0.40\textwidth}
\begin{Keluaran}[]
12 2
\end{Keluaran}
\begin{Catatan}
Hanya bilangan genap yang dijumlahkan: f(6) = 6 + 4 + 2 = 12, f(3) = 2.
\end{Catatan}
\end{column}
\end{columns}
\note{Tanyakan jawaban yang paling sering salah, lalu telusuri bersama.}
\end{frame}

\begin{frame}{Error yang sering muncul minggu ini}
\framesubtitle{Pesan diambil langsung dari g++}
\footnotesize
\begin{tabular}{@{}>{\raggedright\arraybackslash}p{208pt}>{\raggedright\arraybackslash}p{256pt}>{\raggedright\arraybackslash}p{398pt}@{}}
\kepala{Kesalahan} & \kepala{Contoh} & \kepala{Pesan pertama dari g++}\\
Lupa ; sesudah struct & \texttt{struct A \{ int x; \}} & \texttt{error: expected ';' after struct definition}\\
\barisgenap Field tidak ada & \texttt{m.umur pada Mahasiswa} & \texttt{error: 'struct Mahasiswa' has no member named 'umur'}\\
Akses field tanpa indeks & \texttt{rak.judul pada array} & \texttt{error: request for member 'judul' in 'rak', which is of non-class type 'Buku ...}\\
\barisgenap Rekursi tanpa return di satu jalur & \texttt{if (n == 0) return 1; tanpa return lain} & \texttt{warning: control reaches end of non-void function}\\
Rekursi tanpa basis & \texttt{return n * f(n - 1); saja} & \texttt{warning: infinite recursion detected}\\
\garisdasar
\end{tabular}
\vspace{10pt}

{\small Pesan di tabel ini dihasilkan dengan g++ dan opsi -Wall, sama seperti di cek.sh.}
\note{Minta mahasiswa memotret slide ini.}
\end{frame}

\Bagian{4}{Hands-on bertingkat}{45 menit, dari dasar sampai tantangan}
\note{Mahasiswa membuka folder 02 Hands-on/P14 - Rekursi dan Struct - Hands-on.}

\begin{frame}{Peta hands-on}
\framesubtitle{Folder 02 Hands-on/P14 - Rekursi dan Struct - Hands-on}
\small
\begin{tabular}{@{}>{\raggedright\arraybackslash}p{36pt}>{\raggedright\arraybackslash}p{267pt}>{\raggedright\arraybackslash}p{110pt}>{\raggedright\arraybackslash}p{83pt}>{\raggedright\arraybackslash}p{341pt}@{}}
\kepala{No} & \kepala{File} & \kepala{Tingkat} & \kepala{Waktu} & \kepala{Yang dilatih}\\
1 & \texttt{handson1-rekursi.cpp} & Dasar & 10' & kasus basis dan langkah rekursif\\
\barisgenap 2 & \texttt{handson2-struct.cpp} & Menengah & 15' & \texttt{struct}, array of struct, maksimum\\
3 & \texttt{handson3-basis.cpp} & Tantangan & 20' & pohon pemanggilan, perbaiki tiga kesalahan\\
\barisgenap + & \texttt{bonus-hanoi.cpp} & Pengayaan & sisa & rekursi bercabang klasik\\
\garisdasar
\end{tabular}
\vspace{10pt}

Kerjakan berurutan. Cocokkan keluaran dengan folder \texttt{expected/}; latihan yang membaca masukan memakai data di folder \texttt{input/}.\note{Saat berkeliling, minta mahasiswa yang mengerjakan latihan tantangan menjelaskan blok PENJELASAN mereka.}
\end{frame}

\begin{frame}{Cara mengecek dan aturannya}
\framesubtitle{Hands-on}
\begin{columns}[T]
\begin{column}{0.47\textwidth}
\begin{Blok}{Di GDB Online}
Salin isi file latihan, lengkapi TODO, klik Run. Bila program membaca masukan, ketik isi file di \texttt{input/}. Bandingkan dengan \texttt{expected/}.
\end{Blok}
\end{column}
\begin{column}{0.47\textwidth}
\begin{BlokGelap}{Di komputer sendiri}
Jalankan \texttt{bash cek.sh}. Skrip mengompilasi dengan \texttt{-Wall -Werror}, memberi masukan dari \texttt{input/}, lalu mencocokkan keluaran.
\end{BlokGelap}
\end{column}
\end{columns}
\vspace{6pt}
\begin{Catatan}
AI boleh untuk memahami pesan error, tidak untuk meminta solusi utuh. Exit Ticket dikerjakan tanpa bantuan apa pun.
\end{Catatan}
\note{Hands-on tidak dinilai langsung; yang dinilai Exit Ticket dan tugas.}
\end{frame}

\Bagian{5}{Exit Ticket dan tugas}{25 menit terakhir, lalu pekerjaan rumah}
\note{Mulai Exit Ticket tepat waktu.}

\begin{frame}{Exit Ticket 14}
\framesubtitle{Individual, 25 menit, tanpa AI dan internet}
\small
\begin{tabular}{@{}>{\raggedright\arraybackslash}p{60pt}>{\raggedright\arraybackslash}p{560pt}>{\raggedright\arraybackslash}p{80pt}>{\raggedright\arraybackslash}p{150pt}@{}}
\kepala{Soal} & \kepala{Isi} & \kepala{Poin} & \kepala{CPMK}\\
A1 & Tulis fungsi rekursif yang menghitung jumlah 1 sampai n & 25 & CPMK0607\\
\barisgenap A2 & Definisikan struct dan tampilkan data dengan nilai field tertinggi dari array of struct & 25 & CPMK0607\\
B1 & Gambar pohon pemanggilan sebuah fungsi rekursif dan tentukan hasilnya & 20 & CPMK0608\\
\barisgenap B2 & Jelaskan apa yang terjadi bila kasus basis sebuah fungsi dihapus & 20 & CPMK0608\\
B3 & Jelaskan keunggulan array of struct dibandingkan beberapa array sejajar & 10 & CPMK0608\\
\garisdasar
\end{tabular}
\vspace{10pt}

{\small Soal A dinilai dari keluaran yang tepat sama. Soal B dinilai dari ketepatan dan alasan. Pelanggaran aturan membuat nilai Exit Ticket ini 0.}
\note{Soal lengkap dibagikan terpisah dan tidak disimpan di repo publik.}
\end{frame}

\begin{frame}{Tugas 14: Perpustakaan Mini dengan Struct dan Rekursi}
\framesubtitle{Dikumpulkan sebelum Pertemuan 15}
\begin{columns}[T]
\begin{column}{0.47\textwidth}
\begin{Blok}{Bagian A, program, 50 poin}
Buat program perpustakaan mini dengan ketentuan berikut. Semua fungsi pencarian dan agregasi wajib rekursif, bukan perulangan. Ketentuannya di slide berikut.
\end{Blok}
\end{column}
\begin{column}{0.47\textwidth}
\begin{BlokGelap}{Bagian B, penjelasan, 50 poin}
Komentar maksimal 150 kata: gambar tumpukan pemanggilan \texttt{hitungTotalStok} untuk tiga buku contoh, jelaskan kasus basis setiap fungsi rekursif, dan alasan memakai struct dibandingkan array sejajar.
\end{BlokGelap}
\end{column}
\end{columns}
\vspace{8pt}
{\small Data di tugas adalah data kalian sendiri. Rubrik lengkap ada di Modul Belajar 14.}
\note{Ingatkan bahwa penjelasan disalin dari AI mudah dikenali karena tidak menyebut pengalaman mereka sendiri.}
\end{frame}

\begin{frame}{Ketentuan Tugas 14}
\framesubtitle{Perpustakaan Mini dengan Struct dan Rekursi}
\footnotesize
\begin{tabular}{@{}>{\raggedright\arraybackslash}p{87pt}>{\raggedright\arraybackslash}p{788pt}@{}}
\kepala{No} & \kepala{Ketentuan Bagian A}\\
1 & \texttt{struct Buku} berisi kodeBuku, judul, penulis, tahun, dan stok\\
\barisgenap 2 & Array of struct dengan maksimal 10 buku dan validasi masukan\\
3 & Fungsi rekursif \texttt{hitungTotalStok(Buku arr[], int n)}\\
\barisgenap 4 & Fungsi rekursif \texttt{cariBukuTerlama(Buku arr[], int n)} yang mengembalikan indeks buku dengan tahun terkecil\\
5 & Fungsi rekursif \texttt{faktorial(int n)} untuk menghitung banyak cara menyusun n buku di rak\\
\barisgenap 6 & Menu: tambah buku, tampilkan semua, cari buku terlama, total stok, banyak susunan, keluar\\
Bonus & Tambahkan pencarian judul secara rekursif.\\
\garisdasar
\end{tabular}
\note{Bacakan ketentuan satu per satu dan tanyakan apakah ada yang belum jelas.}
\end{frame}

\begin{frame}{Rubrik Tugas 14}
\framesubtitle{Empat level, sama di semua instrumen}
\footnotesize
\begin{tabular}{@{}>{\raggedright\arraybackslash}p{166pt}>{\raggedright\arraybackslash}p{44pt}>{\raggedright\arraybackslash}p{166pt}>{\raggedright\arraybackslash}p{156pt}>{\raggedright\arraybackslash}p{146pt}>{\raggedright\arraybackslash}p{146pt}@{}}
\kepala{Kriteria} & \kepala{Poin} & \kepala{Sangat baik 85--100} & \kepala{Baik 70--84} & \kepala{Cukup 55--69} & \kepala{Kurang <55}\\
A1 Program berjalan & 20 & tanpa error dan peringatan & ada peringatan & perlu satu perbaikan & tidak terkompilasi\\
\barisgenap A2 Struct dan fungsi rekursif & 15 & struct lengkap, tiga fungsi rekursif & satu fungsi tidak rekursif & dua tidak rekursif & tanpa struct\\
A3 Menu dan ketepatan & 15 & semua menu tepat, validasi ada & satu salah & dua salah & sebagian besar salah\\
\barisgenap B1 Tumpukan pemanggilan dan basis & 30 & jejak tepat, basis dijelaskan & satu langkah keliru & tanpa jejak & tidak ada atau disalin\\
B2 Cerita error dan pengujian & 20 & pesan, sebab, perbaikan, uji & pesan tidak dikutip & hanya perbaikan & tidak ada\\
\garisdasar
\end{tabular}
\vspace{8pt}

{\small Nilai kriteria = poin × skor level. Bagian A total 50, Bagian B total 50.}
\note{Rubrik ini sama strukturnya setiap minggu; yang berubah hanya isi kriteria A2, A3, dan B1.}
\end{frame}

\begin{frame}{Sebelum Pertemuan 15}
\framesubtitle{Persiapan}
\begin{columns}[T]
\begin{column}{0.31\textwidth}
\begin{Langkah}{1}{Selesaikan}
Hands-on yang belum lulus \texttt{cek.sh}, terutama latihan tantangan.
\end{Langkah}
\end{column}
\begin{column}{0.31\textwidth}
\begin{Langkah}{2}{Kumpulkan}
Tugas 14 sebelum Pertemuan 15 dimulai.
\end{Langkah}
\end{column}
\begin{column}{0.31\textwidth}
\begin{Langkah}{3}{Baca}
Modul Belajar 14 untuk mengulang, lalu bagian awal modul berikutnya.
\end{Langkah}
\end{column}
\end{columns}
\vspace{10pt}
Pertemuan 15 adalah review dan simulasi UAS untuk materi Pertemuan 9 sampai 14. Bawa catatan dan daftar hal yang masih membingungkan.\note{Tanyakan satu hal yang paling membingungkan hari ini untuk dibuka di review minggu depan.}
\end{frame}

\SlidePenutup{Terima kasih}{Sampai jumpa di Pertemuan 15: review dan simulasi UAS}
\note{Slide penutup.}
\end{document}
"
\documentclass[t]{beamer}
\usetheme{ITERA}
% Mode catatan pembicara: aktif bila \catatan didefinisikan sebelum \input.
\ifdefined\catatan
  \setbeameroption{show only notes}
  \setbeamertemplate{note page}{\vspace*{20pt}\hspace*{4pt}\begin{minipage}{900pt}
    {\ITERAKickerFont\fontsize{11pt}{14pt}\selectfont\color{ITERAGoldText}CATATAN PEMBICARA \ \ |\ \ SLIDE \insertframenumber}\par\vspace{4pt}
    {\ITERADisplay\bfseries\fontsize{22pt}{28pt}\selectfont\color{ITERAStructure}\insertframetitle}\par\vspace{12pt}
    \fontsize{15pt}{23pt}\selectfont\insertnote\end{minipage}}
\fi
\tikzset{fase/.style={draw=none,fill=#1,rounded corners=3pt,minimum width=158pt,minimum height=118pt,text width=146pt,align=center}}

\renewcommand\ITERAmeeting{Pertemuan 14}
\begin{document}
\SlideJudul{IF25-11002 PRAKTIKUM PEMROGRAMAN}{Pertemuan 14\\Rekursi dan Struct}{Fungsi yang memanggil dirinya sendiri, dan tipe data buatan sendiri untuk mengelompokkan data}{Tim Pengajar}{Tahun 2026. Sejajar dengan Pertemuan 14 Algoritma Pemrograman: Rekursi dan Struct.}
\note{Buka dengan menanyakan satu hal dari pertemuan lalu yang masih membingungkan.}

\begin{frame}{Dua mata kuliah, satu minggu}
\framesubtitle{Sebelum mulai}
Topik minggu ini sama di dua mata kuliah, dengan peran yang berbeda.
\vspace{10pt}

\begin{columns}[T]
\begin{column}{0.47\textwidth}
\begin{Blok}{Algoritma: Rekursi dan Struct}
Definisi rekursif, kasus basis, pohon rekursi, serta record atau tipe bentukan sebagai kumpulan field.
\end{Blok}
\end{column}
\begin{column}{0.47\textwidth}
\begin{BlokGelap}{Praktikum: Rekursi dan struct}
Menulis fungsi rekursif dengan kasus basis yang tepat, mendefinisikan \texttt{struct}, dan memproses array of struct.
\end{BlokGelap}
\end{column}
\end{columns}
\note{Tanyakan apa yang sudah dibahas di Algoritma minggu ini. Gunakan jawabannya sebagai jembatan ke praktikum.}
\end{frame}

\begin{frame}{Saat pulang nanti, kalian bisa}
\framesubtitle{Target hari ini}
\begin{columns}[T]
\begin{column}{0.47\textwidth}
\begin{Langkah}{1}{Menerapkan}
Menulis fungsi rekursif dengan kasus basis dan langkah rekursif yang tepat, serta mendefinisikan \texttt{struct} dan memproses array of struct

{\footnotesize\color{ITERAInkMuted}Sub-CPMK 14.1, CPMK0607}
\end{Langkah}
\end{column}
\begin{column}{0.47\textwidth}
\begin{Langkah}{2}{Menjelaskan}
Menelusuri pemanggilan rekursif dengan pohon pemanggilan atau tumpukan panggilan, dan menjelaskan peran kasus basis serta akses anggota struct

{\footnotesize\color{ITERAInkMuted}Sub-CPMK 14.2, CPMK0608}
\end{Langkah}
\end{column}
\end{columns}
\vspace{8pt}
\begin{Catatan}
Dua kemampuan ini dinilai sama berat di Exit Ticket, tugas, UTS, dan UAS.
\end{Catatan}
\note{Bacakan target dengan pelan. Kata kuncinya menerapkan dan menjelaskan.}
\end{frame}

\begin{frame}{Ingat minggu lalu}
\framesubtitle{Review, 5 menit}
\begin{columns}[T]
\begin{column}{0.31\textwidth}
\begin{BlokGelap}{Pertanyaan 1}
Apa beda pass by value dan pass by reference?
\end{BlokGelap}
\end{column}
\begin{column}{0.31\textwidth}
\begin{BlokGelap}{Pertanyaan 2}
Mengapa fungsi non-void harus punya \texttt{return} di semua jalur?
\end{BlokGelap}
\end{column}
\begin{column}{0.31\textwidth}
\begin{BlokGelap}{Pertanyaan 3}
Apakah array di dalam fungsi adalah salinan?
\end{BlokGelap}
\end{column}
\end{columns}
\vspace{8pt}
Jawab di kertas dulu, lalu cocokkan dengan teman sebelah.\note{Pilih dua mahasiswa untuk menjawab. Koreksi singkat, jangan mengajar ulang.}
\end{frame}

\Bagian{1}{Masalah pembuka}{Data buku perpustakaan}
\note{Masalah pembuka 10 menit.}

\begin{frame}[fragile]{Data buku perpustakaan}
\framesubtitle{Masalah minggu ini}
\begin{columns}[T]
\begin{column}{0.55\textwidth}
{\small Perpustakaan prodi mencatat setiap buku dengan kode, judul, penulis, tahun, dan stok. Dengan materi sejauh ini, kalian butuh lima array sejajar dan harus menjaga indeksnya selalu sama.

\vspace{6pt}
Selain itu, beberapa pertanyaan lebih mudah dirumuskan dari ukuran yang lebih kecil, misalnya: total stok n buku adalah stok buku terakhir ditambah total stok n − 1 buku sebelumnya.\par}
\vspace{6pt}
\begin{Catatan}
\textbf{Diskusi 2 menit.} Bagaimana menyimpan lima keterangan satu buku dalam satu variabel? Bagaimana menghitung total stok tanpa perulangan?
\end{Catatan}
\end{column}
\begin{column}{0.41\textwidth}
\begin{Keluaran}[]
[B01] Algoritma (2019) stok 5
[B02] C++ Dasar (2021) stok 8
[B03] Struktur Data (2018) stok 10
Total stok: 23
Buku terlama: Struktur Data
\end{Keluaran}
\end{column}
\end{columns}
\note{Tampung beberapa jawaban tanpa menilai benar salah.}
\end{frame}

\begin{frame}{Dari langkah ke instruksi}
\framesubtitle{Sambungan dengan Algoritma}
\begin{columns}[T]
\begin{column}{0.31\textwidth}
\begin{Langkah}{1}{Kelompokkan}
Satu buku punya beberapa field: kode, judul, penulis, tahun, stok.
\end{Langkah}
\end{column}
\begin{column}{0.31\textwidth}
\begin{Langkah}{2}{Perkecil masalah}
Total n buku = buku terakhir + total n − 1 buku; total 0 buku = 0.
\end{Langkah}
\end{column}
\begin{column}{0.31\textwidth}
\begin{Langkah}{3}{Terjemahkan}
\texttt{struct Buku} dan fungsi rekursif dengan kasus basis.
\end{Langkah}
\end{column}
\end{columns}
\vspace{10pt}
Langkah 1 dan 2 dilatih di Algoritma. Langkah 3 kita kerjakan hari ini dengan C++.\note{Hubungkan eksplisit ke Algoritma minggu ini.}
\end{frame}

\Bagian{2}{Coba dulu, baru dijelaskan}{25 menit eksperimen di GDB Online}
\note{Struggle 25 menit. Berkeliling, jangan menjelaskan materi dulu.}

\begin{frame}{Misi kalian}
\framesubtitle{Struggle, 25 menit}
\begin{columns}[T]
\begin{column}{0.31\textwidth}
\begin{Langkah}{1}{Hitung mundur}
Tulis fungsi yang menampilkan n, n − 1, sampai 1 tanpa \texttt{for} dan \texttt{while}.
\end{Langkah}
\end{column}
\begin{column}{0.31\textwidth}
\begin{Langkah}{2}{Berhenti}
Apa yang terjadi bila fungsi itu tidak punya syarat berhenti?
\end{Langkah}
\end{column}
\begin{column}{0.31\textwidth}
\begin{Langkah}{3}{Kelompokkan}
Simpan nama, NIM, dan IPK satu mahasiswa dalam satu variabel.
\end{Langkah}
\end{column}
\end{columns}
\vspace{8pt}
\begin{columns}[T]
\begin{column}{0.47\textwidth}
\begin{Blok}{Boleh}
Bertanya ke teman, mengubah apa saja, sengaja membuat error.
\end{Blok}
\end{column}
\begin{column}{0.47\textwidth}
\begin{BlokAlert}{Belum dulu}
Mencari jawaban jadi di internet atau AI.
\end{BlokAlert}
\end{column}
\end{columns}
\note{Catat kesalahan yang sering muncul untuk dibahas di debrief.}
\end{frame}

\begin{frame}{Apa yang kalian temukan?}
\framesubtitle{Debrief struggle}
\begin{columns}[T]
\begin{column}{0.31\textwidth}
\begin{BlokGelap}{Pertanyaan 1}
Bagaimana fungsi bisa memanggil dirinya sendiri?
\end{BlokGelap}
\end{column}
\begin{column}{0.31\textwidth}
\begin{BlokGelap}{Pertanyaan 2}
Syarat apa yang membuatnya berhenti?
\end{BlokGelap}
\end{column}
\begin{column}{0.31\textwidth}
\begin{BlokGelap}{Pertanyaan 3}
Apa yang kalian coba untuk menyatukan tiga data mahasiswa?
\end{BlokGelap}
\end{column}
\end{columns}
\vspace{10pt}
Jawaban kalian akan kita uji di bagian berikutnya.\note{Tulis jawaban di papan; buktikan atau patahkan di bagian materi.}
\end{frame}

\Bagian{3}{Memanggil diri sendiri dan tipe buatan}{Kasus basis, pohon pemanggilan, struct, dan array of struct}
\note{Materi 40 menit. Tunjukkan setiap contoh langsung di GDB Online.}

\begin{frame}[fragile]{Rekursi: kasus basis dan langkah rekursif}
\framesubtitle{Faktorial}
\begin{columns}[T]
\begin{column}{0.55\textwidth}
\begin{Kode}[]{faktorial.cpp}
int faktorial(int n) {
    if (n == 0) return 1;
    return n * faktorial(n - 1);
}

int main() {
    cout << faktorial(5) << endl;
    cout << faktorial(0) << endl;
    return 0;
}
\end{Kode}
\end{column}
\begin{column}{0.41\textwidth}
\only<1>{\begin{TebakDulu}
{\ITERAKickerFont\fontsize{10pt}{12pt}\selectfont\color{ITERAAlert}TEBAK DULU}\par\vspace{4pt}
Sebelum dijalankan, tulis keluarannya di kertas.
\end{TebakDulu}
}\only<2>{\KeluaranFile[]{keluaran/P14/k001.txt}
\begin{Catatan}
Kasus basis menghentikan rekursi. Langkah rekursif harus selalu mendekati kasus basis.
\end{Catatan}
}
\end{column}
\end{columns}
\note<1>{Minta mahasiswa menebak keluaran sebelum membuka slide berikutnya. Tanyakan satu atau dua tebakan yang berbeda, lalu buktikan.}
\end{frame}

\begin{frame}{Menelusuri faktorial(4)}
\framesubtitle{Tumpukan pemanggilan}
\small
\begin{tabular}{@{}>{\raggedright\arraybackslash}p{249pt}>{\raggedright\arraybackslash}p{307pt}>{\raggedright\arraybackslash}p{307pt}@{}}
\kepala{Pemanggilan} & \kepala{Menunggu} & \kepala{Nilai kembali}\\
faktorial(4) & 4 × faktorial(3) & 4 × 6 = 24\\
\barisgenap faktorial(3) & 3 × faktorial(2) & 3 × 2 = 6\\
faktorial(2) & 2 × faktorial(1) & 2 × 1 = 2\\
\barisgenap faktorial(1) & 1 × faktorial(0) & 1 × 1 = 1\\
faktorial(0) & kasus basis & 1\\
\garisdasar
\end{tabular}
\vspace{10pt}

Pemanggilan turun sampai kasus basis, lalu nilai kembali naik dari bawah ke atas.
\end{frame}

\begin{frame}[fragile]{Rekursi bercabang dan rekursi pada array}
\framesubtitle{Fibonacci dan jumlah array}
\begin{columns}[T]
\begin{column}{0.60\textwidth}
\begin{Kode}[listing options={style=ITERACodeMini}]{fibarray.cpp}
int fib(int n) {
    if (n <= 1) return n;
    return fib(n - 1) + fib(n - 2);
}
int jumlah(int a[], int n) {
    if (n == 0) return 0;
    return a[n - 1] + jumlah(a, n - 1);
}
int main() {
    int a[4] = {5, 8, 10, 2};
    cout << fib(7) << " " << jumlah(a, 4) << endl;
    return 0;
}
\end{Kode}
\end{column}
\begin{column}{0.36\textwidth}
\only<1>{\begin{TebakDulu}
{\ITERAKickerFont\fontsize{10pt}{12pt}\selectfont\color{ITERAAlert}TEBAK DULU}\par\vspace{4pt}
Sebelum dijalankan, tulis keluarannya di kertas.
\end{TebakDulu}
}\only<2>{\KeluaranFile[listing options={style=ITERAPlainKecil}]{keluaran/P14/k002.txt}
\begin{Catatan}
\texttt{fib} memanggil dirinya dua kali, sehingga pohon pemanggilannya melebar. \texttt{jumlah} memproses elemen terakhir lalu sisanya.
\end{Catatan}
}
\end{column}
\end{columns}
\note<1>{Minta mahasiswa menebak keluaran sebelum membuka slide berikutnya. Tanyakan satu atau dua tebakan yang berbeda, lalu buktikan.}
\end{frame}

\begin{frame}[fragile]{struct: tipe data buatan sendiri}
\framesubtitle{Mengelompokkan field}
\begin{columns}[T]
\begin{column}{0.60\textwidth}
\begin{Kode}[listing options={style=ITERACodeMini}]{struct.cpp}
struct Mahasiswa {
    string nama;
    string nim;
    double ipk;
};
int main() {
    Mahasiswa m;
    m.nama = "Rani";
    m.nim = "126140001";
    m.ipk = 3.72;
    cout << m.nama << " " << m.nim << " " << m.ipk << endl;
    return 0;
}
\end{Kode}
\end{column}
\begin{column}{0.36\textwidth}
\only<1>{\begin{TebakDulu}
{\ITERAKickerFont\fontsize{10pt}{12pt}\selectfont\color{ITERAAlert}TEBAK DULU}\par\vspace{4pt}
Sebelum dijalankan, tulis keluarannya di kertas.
\end{TebakDulu}
}\only<2>{\KeluaranFile[listing options={style=ITERAPlainKecil}]{keluaran/P14/k003.txt}
\begin{Catatan}
Definisi \texttt{struct} diakhiri titik koma. Field diakses dengan titik: \texttt{m.nama}.
\end{Catatan}
}
\end{column}
\end{columns}
\note<1>{Minta mahasiswa menebak keluaran sebelum membuka slide berikutnya. Tanyakan satu atau dua tebakan yang berbeda, lalu buktikan.}
\end{frame}

\begin{frame}[fragile]{Array of struct}
\framesubtitle{Banyak data terstruktur}
\begin{columns}[T]
\begin{column}{0.60\textwidth}
\begin{Kode}[listing options={style=ITERACodeKecil}]{arraystruct.cpp}
struct Buku { string judul; int tahun; };

int main() {
    Buku rak[3] = {{"Algoritma", 2019},
                   {"C++ Dasar", 2021},
                   {"Struktur Data", 2018}};
    int iLama = 0;
    for (int i = 1; i < 3; i++)
        if (rak[i].tahun < rak[iLama].tahun) iLama = i;
    cout << rak[iLama].judul << endl;
    return 0;
}
\end{Kode}
\end{column}
\begin{column}{0.36\textwidth}
\only<1>{\begin{TebakDulu}
{\ITERAKickerFont\fontsize{10pt}{12pt}\selectfont\color{ITERAAlert}TEBAK DULU}\par\vspace{4pt}
Sebelum dijalankan, tulis keluarannya di kertas.
\end{TebakDulu}
}\only<2>{\KeluaranFile[listing options={style=ITERAPlainKecil}]{keluaran/P14/k004.txt}
\begin{Catatan}
\texttt{rak[i].tahun}: indeks dulu untuk memilih buku, lalu titik untuk memilih field.
\end{Catatan}
}
\end{column}
\end{columns}
\note<1>{Minta mahasiswa menebak keluaran sebelum membuka slide berikutnya. Tanyakan satu atau dua tebakan yang berbeda, lalu buktikan.}
\end{frame}

\begin{frame}[fragile]{Tebak keluarannya}
\framesubtitle{Latihan menjelaskan, 3 menit}
\begin{columns}[T]
\begin{column}{0.56\textwidth}
\begin{Kode}[]{tebak.cpp}
int f(int n) {
    if (n <= 0) return 0;
    if (n % 2 == 0) return n + f(n - 1);
    return f(n - 1);
}

int main() {
    cout << f(6) << " " << f(3) << endl;
    return 0;
}
\end{Kode}
\end{column}
\begin{column}{0.40\textwidth}
\begin{Catatan}
Tulis keluarannya di kertas, karakter demi karakter. Jangan dijalankan dulu.
\end{Catatan}
\end{column}
\end{columns}
\note{Contoh soal Bagian B. Beri waktu 3 menit sebelum slide jawaban.}
\end{frame}

\begin{frame}[fragile]{Tebak keluarannya: jawaban}
\framesubtitle{Latihan menjelaskan}
\begin{columns}[T]
\begin{column}{0.56\textwidth}
\begin{Kode}[]{tebak.cpp}
int f(int n) {
    if (n <= 0) return 0;
    if (n % 2 == 0) return n + f(n - 1);
    return f(n - 1);
}

int main() {
    cout << f(6) << " " << f(3) << endl;
    return 0;
}
\end{Kode}
\end{column}
\begin{column}{0.40\textwidth}
\begin{Keluaran}[]
12 2
\end{Keluaran}
\begin{Catatan}
Hanya bilangan genap yang dijumlahkan: f(6) = 6 + 4 + 2 = 12, f(3) = 2.
\end{Catatan}
\end{column}
\end{columns}
\note{Tanyakan jawaban yang paling sering salah, lalu telusuri bersama.}
\end{frame}

\begin{frame}{Error yang sering muncul minggu ini}
\framesubtitle{Pesan diambil langsung dari g++}
\footnotesize
\begin{tabular}{@{}>{\raggedright\arraybackslash}p{208pt}>{\raggedright\arraybackslash}p{256pt}>{\raggedright\arraybackslash}p{398pt}@{}}
\kepala{Kesalahan} & \kepala{Contoh} & \kepala{Pesan pertama dari g++}\\
Lupa ; sesudah struct & \texttt{struct A \{ int x; \}} & \texttt{error: expected ';' after struct definition}\\
\barisgenap Field tidak ada & \texttt{m.umur pada Mahasiswa} & \texttt{error: 'struct Mahasiswa' has no member named 'umur'}\\
Akses field tanpa indeks & \texttt{rak.judul pada array} & \texttt{error: request for member 'judul' in 'rak', which is of non-class type 'Buku ...}\\
\barisgenap Rekursi tanpa return di satu jalur & \texttt{if (n == 0) return 1; tanpa return lain} & \texttt{warning: control reaches end of non-void function}\\
Rekursi tanpa basis & \texttt{return n * f(n - 1); saja} & \texttt{warning: infinite recursion detected}\\
\garisdasar
\end{tabular}
\vspace{10pt}

{\small Pesan di tabel ini dihasilkan dengan g++ dan opsi -Wall, sama seperti di cek.sh.}
\note{Minta mahasiswa memotret slide ini.}
\end{frame}

\Bagian{4}{Hands-on bertingkat}{45 menit, dari dasar sampai tantangan}
\note{Mahasiswa membuka folder 02 Hands-on/P14 - Rekursi dan Struct - Hands-on.}

\begin{frame}{Peta hands-on}
\framesubtitle{Folder 02 Hands-on/P14 - Rekursi dan Struct - Hands-on}
\small
\begin{tabular}{@{}>{\raggedright\arraybackslash}p{36pt}>{\raggedright\arraybackslash}p{267pt}>{\raggedright\arraybackslash}p{110pt}>{\raggedright\arraybackslash}p{83pt}>{\raggedright\arraybackslash}p{341pt}@{}}
\kepala{No} & \kepala{File} & \kepala{Tingkat} & \kepala{Waktu} & \kepala{Yang dilatih}\\
1 & \texttt{handson1-rekursi.cpp} & Dasar & 10' & kasus basis dan langkah rekursif\\
\barisgenap 2 & \texttt{handson2-struct.cpp} & Menengah & 15' & \texttt{struct}, array of struct, maksimum\\
3 & \texttt{handson3-basis.cpp} & Tantangan & 20' & pohon pemanggilan, perbaiki tiga kesalahan\\
\barisgenap + & \texttt{bonus-hanoi.cpp} & Pengayaan & sisa & rekursi bercabang klasik\\
\garisdasar
\end{tabular}
\vspace{10pt}

Kerjakan berurutan. Cocokkan keluaran dengan folder \texttt{expected/}; latihan yang membaca masukan memakai data di folder \texttt{input/}.\note{Saat berkeliling, minta mahasiswa yang mengerjakan latihan tantangan menjelaskan blok PENJELASAN mereka.}
\end{frame}

\begin{frame}{Cara mengecek dan aturannya}
\framesubtitle{Hands-on}
\begin{columns}[T]
\begin{column}{0.47\textwidth}
\begin{Blok}{Di GDB Online}
Salin isi file latihan, lengkapi TODO, klik Run. Bila program membaca masukan, ketik isi file di \texttt{input/}. Bandingkan dengan \texttt{expected/}.
\end{Blok}
\end{column}
\begin{column}{0.47\textwidth}
\begin{BlokGelap}{Di komputer sendiri}
Jalankan \texttt{bash cek.sh}. Skrip mengompilasi dengan \texttt{-Wall -Werror}, memberi masukan dari \texttt{input/}, lalu mencocokkan keluaran.
\end{BlokGelap}
\end{column}
\end{columns}
\vspace{6pt}
\begin{Catatan}
AI boleh untuk memahami pesan error, tidak untuk meminta solusi utuh. Exit Ticket dikerjakan tanpa bantuan apa pun.
\end{Catatan}
\note{Hands-on tidak dinilai langsung; yang dinilai Exit Ticket dan tugas.}
\end{frame}

\Bagian{5}{Exit Ticket dan tugas}{25 menit terakhir, lalu pekerjaan rumah}
\note{Mulai Exit Ticket tepat waktu.}

\begin{frame}{Exit Ticket 14}
\framesubtitle{Individual, 25 menit, tanpa AI dan internet}
\small
\begin{tabular}{@{}>{\raggedright\arraybackslash}p{60pt}>{\raggedright\arraybackslash}p{560pt}>{\raggedright\arraybackslash}p{80pt}>{\raggedright\arraybackslash}p{150pt}@{}}
\kepala{Soal} & \kepala{Isi} & \kepala{Poin} & \kepala{CPMK}\\
A1 & Tulis fungsi rekursif yang menghitung jumlah 1 sampai n & 25 & CPMK0607\\
\barisgenap A2 & Definisikan struct dan tampilkan data dengan nilai field tertinggi dari array of struct & 25 & CPMK0607\\
B1 & Gambar pohon pemanggilan sebuah fungsi rekursif dan tentukan hasilnya & 20 & CPMK0608\\
\barisgenap B2 & Jelaskan apa yang terjadi bila kasus basis sebuah fungsi dihapus & 20 & CPMK0608\\
B3 & Jelaskan keunggulan array of struct dibandingkan beberapa array sejajar & 10 & CPMK0608\\
\garisdasar
\end{tabular}
\vspace{10pt}

{\small Soal A dinilai dari keluaran yang tepat sama. Soal B dinilai dari ketepatan dan alasan. Pelanggaran aturan membuat nilai Exit Ticket ini 0.}
\note{Soal lengkap dibagikan terpisah dan tidak disimpan di repo publik.}
\end{frame}

\begin{frame}{Tugas 14: Perpustakaan Mini dengan Struct dan Rekursi}
\framesubtitle{Dikumpulkan sebelum Pertemuan 15}
\begin{columns}[T]
\begin{column}{0.47\textwidth}
\begin{Blok}{Bagian A, program, 50 poin}
Buat program perpustakaan mini dengan ketentuan berikut. Semua fungsi pencarian dan agregasi wajib rekursif, bukan perulangan. Ketentuannya di slide berikut.
\end{Blok}
\end{column}
\begin{column}{0.47\textwidth}
\begin{BlokGelap}{Bagian B, penjelasan, 50 poin}
Komentar maksimal 150 kata: gambar tumpukan pemanggilan \texttt{hitungTotalStok} untuk tiga buku contoh, jelaskan kasus basis setiap fungsi rekursif, dan alasan memakai struct dibandingkan array sejajar.
\end{BlokGelap}
\end{column}
\end{columns}
\vspace{8pt}
{\small Data di tugas adalah data kalian sendiri. Rubrik lengkap ada di Modul Belajar 14.}
\note{Ingatkan bahwa penjelasan disalin dari AI mudah dikenali karena tidak menyebut pengalaman mereka sendiri.}
\end{frame}

\begin{frame}{Ketentuan Tugas 14}
\framesubtitle{Perpustakaan Mini dengan Struct dan Rekursi}
\footnotesize
\begin{tabular}{@{}>{\raggedright\arraybackslash}p{87pt}>{\raggedright\arraybackslash}p{788pt}@{}}
\kepala{No} & \kepala{Ketentuan Bagian A}\\
1 & \texttt{struct Buku} berisi kodeBuku, judul, penulis, tahun, dan stok\\
\barisgenap 2 & Array of struct dengan maksimal 10 buku dan validasi masukan\\
3 & Fungsi rekursif \texttt{hitungTotalStok(Buku arr[], int n)}\\
\barisgenap 4 & Fungsi rekursif \texttt{cariBukuTerlama(Buku arr[], int n)} yang mengembalikan indeks buku dengan tahun terkecil\\
5 & Fungsi rekursif \texttt{faktorial(int n)} untuk menghitung banyak cara menyusun n buku di rak\\
\barisgenap 6 & Menu: tambah buku, tampilkan semua, cari buku terlama, total stok, banyak susunan, keluar\\
Bonus & Tambahkan pencarian judul secara rekursif.\\
\garisdasar
\end{tabular}
\note{Bacakan ketentuan satu per satu dan tanyakan apakah ada yang belum jelas.}
\end{frame}

\begin{frame}{Rubrik Tugas 14}
\framesubtitle{Empat level, sama di semua instrumen}
\footnotesize
\begin{tabular}{@{}>{\raggedright\arraybackslash}p{166pt}>{\raggedright\arraybackslash}p{44pt}>{\raggedright\arraybackslash}p{166pt}>{\raggedright\arraybackslash}p{156pt}>{\raggedright\arraybackslash}p{146pt}>{\raggedright\arraybackslash}p{146pt}@{}}
\kepala{Kriteria} & \kepala{Poin} & \kepala{Sangat baik 85--100} & \kepala{Baik 70--84} & \kepala{Cukup 55--69} & \kepala{Kurang <55}\\
A1 Program berjalan & 20 & tanpa error dan peringatan & ada peringatan & perlu satu perbaikan & tidak terkompilasi\\
\barisgenap A2 Struct dan fungsi rekursif & 15 & struct lengkap, tiga fungsi rekursif & satu fungsi tidak rekursif & dua tidak rekursif & tanpa struct\\
A3 Menu dan ketepatan & 15 & semua menu tepat, validasi ada & satu salah & dua salah & sebagian besar salah\\
\barisgenap B1 Tumpukan pemanggilan dan basis & 30 & jejak tepat, basis dijelaskan & satu langkah keliru & tanpa jejak & tidak ada atau disalin\\
B2 Cerita error dan pengujian & 20 & pesan, sebab, perbaikan, uji & pesan tidak dikutip & hanya perbaikan & tidak ada\\
\garisdasar
\end{tabular}
\vspace{8pt}

{\small Nilai kriteria = poin × skor level. Bagian A total 50, Bagian B total 50.}
\note{Rubrik ini sama strukturnya setiap minggu; yang berubah hanya isi kriteria A2, A3, dan B1.}
\end{frame}

\begin{frame}{Sebelum Pertemuan 15}
\framesubtitle{Persiapan}
\begin{columns}[T]
\begin{column}{0.31\textwidth}
\begin{Langkah}{1}{Selesaikan}
Hands-on yang belum lulus \texttt{cek.sh}, terutama latihan tantangan.
\end{Langkah}
\end{column}
\begin{column}{0.31\textwidth}
\begin{Langkah}{2}{Kumpulkan}
Tugas 14 sebelum Pertemuan 15 dimulai.
\end{Langkah}
\end{column}
\begin{column}{0.31\textwidth}
\begin{Langkah}{3}{Baca}
Modul Belajar 14 untuk mengulang, lalu bagian awal modul berikutnya.
\end{Langkah}
\end{column}
\end{columns}
\vspace{10pt}
Pertemuan 15 adalah review dan simulasi UAS untuk materi Pertemuan 9 sampai 14. Bawa catatan dan daftar hal yang masih membingungkan.\note{Tanyakan satu hal yang paling membingungkan hari ini untuk dibuka di review minggu depan.}
\end{frame}

\SlidePenutup{Terima kasih}{Sampai jumpa di Pertemuan 15: review dan simulasi UAS}
\note{Slide penutup.}
\end{document}
